\section{Detailed provenance for every inventory record}
All 137 records are included, not only the active or sensitive subset. The full original JSON and CSV additionally preserve all raw metadata and complete code location lists. Repeated generic caveats are not evidence of an independently measured constraint. Values, units, status and the following explanatory fields are transcribed from the audit without turning logical domains into physiological uncertainty intervals.
\subsubsection*{Record 1: temperature K}
{\small \nolinkurl{constants.temperature_K}\par

Value: \texttt{310.15}; units: K; active: Yes; classification: effective\_assumption.\par

Default: 310.15; differs from default: False.\par

\textbf{Equation role.} Thermal voltage RT/F, Nernst potentials, and conversion between current and ionic amount flux.\par

\textbf{Constraint.} 37 C reference temperature; protocol-specific temperature must replace it\par

\textbf{Limit of justification.} An accepted physiological point does not identify this setting or supply a numerical uncertainty interval. No knockout fit or new measurement is performed in this audit.\par

\textbf{Logical domain.} Positive for the retained law; no finite upper physiological bound established.\par

\textbf{Uncertainty status.} genuinely\_unknown\_for\_physiological\_transfer\par

\textbf{Dependence of claims.} Numerical state, current and secretion claims are conditional on this active setting; exact stoichiometric identities are independent of its numerical value.\par

\textbf{Discriminating measurement.} Measure this quantity or its identifiable combination in the same preparation and stimulus protocol, with independent WT data; determine physiological ranges before making global robustness claims.\par

\textbf{Source.} \nolinkurl{src/modern_full_model/membranes.py:50}. No independently verified primary measurement attached to this numerical setting; model lineage is the source.\par}

\subsubsection*{Record 2: gas constant J mol K}
{\small \nolinkurl{constants.gas_constant_J_mol_K}\par

Value: \texttt{8.31446261815324}; units: J mol\textasciicircum{}-1 K\textasciicircum{}-1; active: Yes; classification: physical\_constant.\par

Default: 8.31446261815324; differs from default: False.\par

\textbf{Equation role.} Thermal voltage RT/F, Nernst potentials, and conversion between current and ionic amount flux.\par

\textbf{Constraint.} CODATA exact/standard physical constant\par

\textbf{Limit of justification.} No physiological fitting or salivary uncertainty interval applies to a standard physical constant.\par

\textbf{Logical domain.} Positive for the retained law; no finite upper physiological bound established.\par

\textbf{Uncertainty status.} fixed\_standard\_constant\_at\_stored\_precision\par

\textbf{Dependence of claims.} Numerical state, current and secretion claims are conditional on this active setting; exact stoichiometric identities are independent of its numerical value.\par

\textbf{Discriminating measurement.} No physiological experiment required; retain standard value and unit convention.\par

\textbf{Source.} \nolinkurl{src/modern_full_model/membranes.py:50}. Standard physical constant, rounded to the stored precision; not a physiological estimate.\par}

\subsubsection*{Record 3: faraday C mol}
{\small \nolinkurl{constants.faraday_C_mol}\par

Value: \texttt{96485.33212000001}; units: C mol\textasciicircum{}-1; active: Yes; classification: physical\_constant.\par

Default: 96485.33212; differs from default: False.\par

\textbf{Equation role.} Thermal voltage RT/F, Nernst potentials, and conversion between current and ionic amount flux.\par

\textbf{Constraint.} Faraday conversion between charge current and molar flux\par

\textbf{Limit of justification.} No physiological fitting or salivary uncertainty interval applies to a standard physical constant.\par

\textbf{Logical domain.} Positive for the retained law; no finite upper physiological bound established.\par

\textbf{Uncertainty status.} fixed\_standard\_constant\_at\_stored\_precision\par

\textbf{Dependence of claims.} Numerical state, current and secretion claims are conditional on this active setting; exact stoichiometric identities are independent of its numerical value.\par

\textbf{Discriminating measurement.} No physiological experiment required; retain standard value and unit convention.\par

\textbf{Source.} \nolinkurl{src/modern_full_model/membranes.py:50}. Standard physical constant, rounded to the stored precision; not a physiological estimate.\par}

\subsubsection*{Record 4: carbon pka1}
{\small \nolinkurl{acid_base.carbon_pka1}\par

Value: \texttt{6.1}; units: dimensionless (p-scale); active: Yes; classification: inherited\_model\_setting.\par

Default: 6.1; differs from default: False.\par

\textbf{Equation role.} Algebraic TA = carbonate alkalinity + finite buffer base + OH - H; pH is solved from TIC, TA and volume.\par

\textbf{Constraint.} effective CO2/HCO3 pKa; temperature/context refinement required Source metadata identifies a model lineage; independent experimental measurement is not established here.\par

\textbf{Limit of justification.} An accepted physiological point does not identify this setting or supply a numerical uncertainty interval. No knockout fit or new measurement is performed in this audit.\par

\textbf{Logical domain.} Finite p scale dissociation constant; environmental applicability must be measured.\par

\textbf{Uncertainty status.} genuinely\_unknown\_for\_physiological\_transfer\par

\textbf{Dependence of claims.} Numerical state, current and secretion claims are conditional on this active setting; exact stoichiometric identities are independent of its numerical value.\par

\textbf{Discriminating measurement.} Measure this quantity or its identifiable combination in the same preparation and stimulus protocol, with independent WT data; determine physiological ranges before making global robustness claims.\par

\textbf{Source.} \nolinkurl{src/modern_full_model/acid_base.py:93}. No independently verified primary measurement attached to this numerical setting; model lineage is the source.\par}

\subsubsection*{Record 5: carbon pka2}
{\small \nolinkurl{acid_base.carbon_pka2}\par

Value: \texttt{10.3}; units: dimensionless (p-scale); active: Yes; classification: inherited\_model\_setting.\par

Default: 10.3; differs from default: False.\par

\textbf{Equation role.} Algebraic TA = carbonate alkalinity + finite buffer base + OH - H; pH is solved from TIC, TA and volume.\par

\textbf{Constraint.} effective HCO3/CO3 pKa; carbonate retained for charge closure Source metadata identifies a model lineage; independent experimental measurement is not established here.\par

\textbf{Limit of justification.} An accepted physiological point does not identify this setting or supply a numerical uncertainty interval. No knockout fit or new measurement is performed in this audit.\par

\textbf{Logical domain.} Finite p scale dissociation constant; environmental applicability must be measured.\par

\textbf{Uncertainty status.} genuinely\_unknown\_for\_physiological\_transfer\par

\textbf{Dependence of claims.} Numerical state, current and secretion claims are conditional on this active setting; exact stoichiometric identities are independent of its numerical value.\par

\textbf{Discriminating measurement.} Measure this quantity or its identifiable combination in the same preparation and stimulus protocol, with independent WT data; determine physiological ranges before making global robustness claims.\par

\textbf{Source.} \nolinkurl{src/modern_full_model/acid_base.py:93}. No independently verified primary measurement attached to this numerical setting; model lineage is the source.\par}

\subsubsection*{Record 6: water pkw}
{\small \nolinkurl{acid_base.water_pkw}\par

Value: \texttt{14}; units: dimensionless (p-scale); active: Yes; classification: inherited\_model\_setting.\par

Default: 14.0; differs from default: False.\par

\textbf{Equation role.} Algebraic TA = carbonate alkalinity + finite buffer base + OH - H; pH is solved from TIC, TA and volume.\par

\textbf{Constraint.} effective water ionization value; temperature refinement required Source metadata identifies a model lineage; independent experimental measurement is not established here.\par

\textbf{Limit of justification.} An accepted physiological point does not identify this setting or supply a numerical uncertainty interval. No knockout fit or new measurement is performed in this audit.\par

\textbf{Logical domain.} Finite p scale dissociation constant; environmental applicability must be measured.\par

\textbf{Uncertainty status.} genuinely\_unknown\_for\_physiological\_transfer\par

\textbf{Dependence of claims.} Numerical state, current and secretion claims are conditional on this active setting; exact stoichiometric identities are independent of its numerical value.\par

\textbf{Discriminating measurement.} Measure this quantity or its identifiable combination in the same preparation and stimulus protocol, with independent WT data; determine physiological ranges before making global robustness claims.\par

\textbf{Source.} \nolinkurl{src/modern_full_model/acid_base.py:93}. No independently verified primary measurement attached to this numerical setting; model lineage is the source.\par}

\subsubsection*{Record 7: cell buffer pka}
{\small \nolinkurl{acid_base.cell_buffer_pka}\par

Value: \texttt{7}; units: dimensionless (p-scale); active: Yes; classification: effective\_assumption.\par

Default: 7.0; differs from default: False.\par

\textbf{Equation role.} Algebraic TA = carbonate alkalinity + finite buffer base + OH - H; pH is solved from TIC, TA and volume.\par

\textbf{Constraint.} single effective non-carbonate buffer site; not identified\par

\textbf{Limit of justification.} An accepted physiological point does not identify this setting or supply a numerical uncertainty interval. No knockout fit or new measurement is performed in this audit.\par

\textbf{Logical domain.} Finite p scale dissociation constant; environmental applicability must be measured.\par

\textbf{Uncertainty status.} genuinely\_unknown\_for\_physiological\_transfer\par

\textbf{Dependence of claims.} Numerical state, current and secretion claims are conditional on this active setting; exact stoichiometric identities are independent of its numerical value.\par

\textbf{Discriminating measurement.} Measure this quantity or its identifiable combination in the same preparation and stimulus protocol, with independent WT data; determine physiological ranges before making global robustness claims.\par

\textbf{Source.} \nolinkurl{src/modern_full_model/acid_base.py:93}. No independently verified primary measurement attached to this numerical setting; model lineage is the source.\par}

\subsubsection*{Record 8: lumen buffer pka}
{\small \nolinkurl{acid_base.lumen_buffer_pka}\par

Value: \texttt{7}; units: dimensionless (p-scale); active: Yes; classification: effective\_assumption.\par

Default: 7.0; differs from default: False.\par

\textbf{Equation role.} Algebraic TA = carbonate alkalinity + finite buffer base + OH - H; pH is solved from TIC, TA and volume.\par

\textbf{Constraint.} irrelevant when luminal buffer amount is zero\par

\textbf{Limit of justification.} An accepted physiological point does not identify this setting or supply a numerical uncertainty interval. No knockout fit or new measurement is performed in this audit.\par

\textbf{Logical domain.} Finite p scale dissociation constant; environmental applicability must be measured.\par

\textbf{Uncertainty status.} genuinely\_unknown\_for\_physiological\_transfer\par

\textbf{Dependence of claims.} Numerical state, current and secretion claims are conditional on this active setting; exact stoichiometric identities are independent of its numerical value.\par

\textbf{Discriminating measurement.} Measure this quantity or its identifiable combination in the same preparation and stimulus protocol, with independent WT data; determine physiological ranges before making global robustness claims.\par

\textbf{Source.} \nolinkurl{src/modern_full_model/acid_base.py:93}. No independently verified primary measurement attached to this numerical setting; model lineage is the source.\par}

\subsubsection*{Record 9: ph lower}
{\small \nolinkurl{acid_base.ph_lower}\par

Value: \texttt{3}; units: dimensionless (p-scale); active: Yes; classification: numerical\_bracket.\par

Default: 3.0; differs from default: False.\par

\textbf{Equation role.} Algebraic TA = carbonate alkalinity + finite buffer base + OH - H; pH is solved from TIC, TA and volume.\par

\textbf{Constraint.} Numerical pH inversion bracket, not an experimental physiological limit.\par

\textbf{Limit of justification.} An accepted physiological point does not identify this setting or supply a numerical uncertainty interval. No knockout fit or new measurement is performed in this audit.\par

\textbf{Logical domain.} Ordered numerical inversion endpoints [3,11]; not the inherited physiological pH gate [6.6,7.3].\par

\textbf{Uncertainty status.} not\_a\_measured\_uncertainty\_parameter\par

\textbf{Dependence of claims.} Numerical state, current and secretion claims are conditional on this active setting; exact stoichiometric identities are independent of its numerical value.\par

\textbf{Discriminating measurement.} Verify numerical convergence and successful closure on the stated domain; this is not an abundance measurement.\par

\textbf{Source.} \nolinkurl{src/modern_full_model/acid_base.py:93}. No independently verified primary measurement attached to this numerical setting; model lineage is the source.\par}

\subsubsection*{Record 10: ph upper}
{\small \nolinkurl{acid_base.ph_upper}\par

Value: \texttt{11}; units: dimensionless (p-scale); active: Yes; classification: numerical\_bracket.\par

Default: 11.0; differs from default: False.\par

\textbf{Equation role.} Algebraic TA = carbonate alkalinity + finite buffer base + OH - H; pH is solved from TIC, TA and volume.\par

\textbf{Constraint.} Numerical pH inversion bracket, not an experimental physiological limit.\par

\textbf{Limit of justification.} An accepted physiological point does not identify this setting or supply a numerical uncertainty interval. No knockout fit or new measurement is performed in this audit.\par

\textbf{Logical domain.} Ordered numerical inversion endpoints [3,11]; not the inherited physiological pH gate [6.6,7.3].\par

\textbf{Uncertainty status.} not\_a\_measured\_uncertainty\_parameter\par

\textbf{Dependence of claims.} Numerical state, current and secretion claims are conditional on this active setting; exact stoichiometric identities are independent of its numerical value.\par

\textbf{Discriminating measurement.} Verify numerical convergence and successful closure on the stated domain; this is not an abundance measurement.\par

\textbf{Source.} \nolinkurl{src/modern_full_model/acid_base.py:93}. No independently verified primary measurement attached to this numerical setting; model lineage is the source.\par}

\subsubsection*{Record 11: na mM}
{\small \nolinkurl{bath.na_mM}\par

Value: \texttt{145}; units: mM; active: Yes; classification: effective\_assumption.\par

Default: 145.0; differs from default: False.\par

\textbf{Equation role.} Fixed reservoir composition determines bath pH species, osmolarity, transporter affinities and reversal potentials.\par

\textbf{Constraint.} representative open-bath value; protocol ledger must replace it\par

\textbf{Limit of justification.} An accepted physiological point does not identify this setting or supply a numerical uncertainty interval. No knockout fit or new measurement is performed in this audit.\par

\textbf{Logical domain.} Positive for the retained law; no finite upper physiological bound established.\par

\textbf{Uncertainty status.} genuinely\_unknown\_for\_physiological\_transfer\par

\textbf{Dependence of claims.} Numerical state, current and secretion claims are conditional on this active setting; exact stoichiometric identities are independent of its numerical value.\par

\textbf{Discriminating measurement.} Measure this quantity or its identifiable combination in the same preparation and stimulus protocol, with independent WT data; determine physiological ranges before making global robustness claims.\par

\textbf{Source.} \nolinkurl{src/modern_full_model/model.py:374}. No independently verified primary measurement attached to this numerical setting; model lineage is the source.\par}

\subsubsection*{Record 12: k mM}
{\small \nolinkurl{bath.k_mM}\par

Value: \texttt{5}; units: mM; active: Yes; classification: effective\_assumption.\par

Default: 5.0; differs from default: False.\par

\textbf{Equation role.} Fixed reservoir composition determines bath pH species, osmolarity, transporter affinities and reversal potentials.\par

\textbf{Constraint.} representative open-bath value; protocol ledger must replace it\par

\textbf{Limit of justification.} An accepted physiological point does not identify this setting or supply a numerical uncertainty interval. No knockout fit or new measurement is performed in this audit.\par

\textbf{Logical domain.} Positive for the retained law; no finite upper physiological bound established.\par

\textbf{Uncertainty status.} genuinely\_unknown\_for\_physiological\_transfer\par

\textbf{Dependence of claims.} Numerical state, current and secretion claims are conditional on this active setting; exact stoichiometric identities are independent of its numerical value.\par

\textbf{Discriminating measurement.} Measure this quantity or its identifiable combination in the same preparation and stimulus protocol, with independent WT data; determine physiological ranges before making global robustness claims.\par

\textbf{Source.} \nolinkurl{src/modern_full_model/model.py:374}. No independently verified primary measurement attached to this numerical setting; model lineage is the source.\par}

\subsubsection*{Record 13: cl mM}
{\small \nolinkurl{bath.cl_mM}\par

Value: \texttt{126.1615928036609}; units: mM; active: Yes; classification: reference\_constraint\_derived.\par

Default: 126.16159280366088; differs from default: False.\par

\textbf{Equation role.} Fixed reservoir composition determines bath pH species, osmolarity, transporter affinities and reversal potentials.\par

\textbf{Constraint.} Bath electroneutrality using the adopted bath sodium, potassium and carbonate speciation.\par

\textbf{Limit of justification.} An accepted physiological point does not identify this setting or supply a numerical uncertainty interval. No knockout fit or new measurement is performed in this audit.\par

\textbf{Logical domain.} Positive for the retained law; no finite upper physiological bound established.\par

\textbf{Uncertainty status.} conditional\_derivation\_no\_experimental\_interval\par

\textbf{Dependence of claims.} Numerical state, current and secretion claims are conditional on this active setting; exact stoichiometric identities are independent of its numerical value.\par

\textbf{Discriminating measurement.} Measure this quantity or its identifiable combination in the same preparation and stimulus protocol, with independent WT data; determine physiological ranges before making global robustness claims.\par

\textbf{Source.} \nolinkurl{src/modern_full_model/model.py:374}. No independently verified primary measurement attached to this numerical setting; model lineage is the source.\par}

\subsubsection*{Record 14: tic mM}
{\small \nolinkurl{bath.tic_mM}\par

Value: \texttt{25}; units: mM; active: Yes; classification: effective\_assumption.\par

Default: 25.0; differs from default: False.\par

\textbf{Equation role.} Fixed reservoir composition determines bath pH species, osmolarity, transporter affinities and reversal potentials.\par

\textbf{Constraint.} open-bath total inorganic carbon starting value\par

\textbf{Limit of justification.} An accepted physiological point does not identify this setting or supply a numerical uncertainty interval. No knockout fit or new measurement is performed in this audit.\par

\textbf{Logical domain.} Positive for the retained law; no finite upper physiological bound established.\par

\textbf{Uncertainty status.} genuinely\_unknown\_for\_physiological\_transfer\par

\textbf{Dependence of claims.} Numerical state, current and secretion claims are conditional on this active setting; exact stoichiometric identities are independent of its numerical value.\par

\textbf{Discriminating measurement.} Measure this quantity or its identifiable combination in the same preparation and stimulus protocol, with independent WT data; determine physiological ranges before making global robustness claims.\par

\textbf{Source.} \nolinkurl{src/modern_full_model/model.py:374}. No independently verified primary measurement attached to this numerical setting; model lineage is the source.\par}

\subsubsection*{Record 15: ph}
{\small \nolinkurl{bath.ph}\par

Value: \texttt{7.4}; units: dimensionless (p-scale); active: Yes; classification: effective\_assumption.\par

Default: 7.4; differs from default: False.\par

\textbf{Equation role.} Fixed reservoir composition determines bath pH species, osmolarity, transporter affinities and reversal potentials.\par

\textbf{Constraint.} representative bicarbonate-buffered bath value\par

\textbf{Limit of justification.} An accepted physiological point does not identify this setting or supply a numerical uncertainty interval. No knockout fit or new measurement is performed in this audit.\par

\textbf{Logical domain.} Positive for the retained law; no finite upper physiological bound established.\par

\textbf{Uncertainty status.} genuinely\_unknown\_for\_physiological\_transfer\par

\textbf{Dependence of claims.} Numerical state, current and secretion claims are conditional on this active setting; exact stoichiometric identities are independent of its numerical value.\par

\textbf{Discriminating measurement.} Measure this quantity or its identifiable combination in the same preparation and stimulus protocol, with independent WT data; determine physiological ranges before making global robustness claims.\par

\textbf{Source.} \nolinkurl{src/modern_full_model/model.py:374}. No independently verified primary measurement attached to this numerical setting; model lineage is the source.\par}

\subsubsection*{Record 16: buffer total mM}
{\small \nolinkurl{bath.buffer_total_mM}\par

Value: \texttt{0}; units: mM sites; active: Yes; classification: effective\_assumption.\par

Default: 0.0; differs from default: False.\par

\textbf{Equation role.} Fixed reservoir composition determines bath pH species, osmolarity, transporter affinities and reversal potentials.\par

\textbf{Constraint.} bath non-carbonate buffer omitted from exchanged species An exact zero is an explicit structural omission, not a small positive fitted value.\par

\textbf{Limit of justification.} An accepted physiological point does not identify this setting or supply a numerical uncertainty interval. No knockout fit or new measurement is performed in this audit.\par

\textbf{Logical domain.} Exact zero records a declared omitted contribution or protocol normalisation; it is not a confidence bound.\par

\textbf{Uncertainty status.} genuinely\_unknown\_for\_physiological\_transfer\par

\textbf{Dependence of claims.} Numerical state, current and secretion claims are conditional on this active setting; exact stoichiometric identities are independent of its numerical value.\par

\textbf{Discriminating measurement.} Measure this quantity or its identifiable combination in the same preparation and stimulus protocol, with independent WT data; determine physiological ranges before making global robustness claims.\par

\textbf{Source.} \nolinkurl{src/modern_full_model/model.py:374}. No independently verified primary measurement attached to this numerical setting; model lineage is the source.\par}

\subsubsection*{Record 17: untracked osmolyte mM}
{\small \nolinkurl{bath.untracked_osmolyte_mM}\par

Value: \texttt{0}; units: mOsm L\textasciicircum{}-1; active: Yes; classification: effective\_assumption.\par

Default: 0.0; differs from default: False.\par

\textbf{Equation role.} Fixed reservoir composition determines bath pH species, osmolarity, transporter affinities and reversal potentials.\par

\textbf{Constraint.} explicit correction only if protocol osmolarity requires it An exact zero is an explicit structural omission, not a small positive fitted value.\par

\textbf{Limit of justification.} An accepted physiological point does not identify this setting or supply a numerical uncertainty interval. No knockout fit or new measurement is performed in this audit.\par

\textbf{Logical domain.} Exact zero records a declared omitted contribution or protocol normalisation; it is not a confidence bound.\par

\textbf{Uncertainty status.} genuinely\_unknown\_for\_physiological\_transfer\par

\textbf{Dependence of claims.} Numerical state, current and secretion claims are conditional on this active setting; exact stoichiometric identities are independent of its numerical value.\par

\textbf{Discriminating measurement.} Measure this quantity or its identifiable combination in the same preparation and stimulus protocol, with independent WT data; determine physiological ranges before making global robustness claims.\par

\textbf{Source.} \nolinkurl{src/modern_full_model/model.py:374}. No independently verified primary measurement attached to this numerical setting; model lineage is the source.\par}

\subsubsection*{Record 18: cell buffer total fmol}
{\small \nolinkurl{geometry.cell_buffer_total_fmol}\par

Value: \texttt{30}; units: fmol sites; active: Yes; classification: effective\_assumption.\par

Default: 30.0; differs from default: False.\par

\textbf{Equation role.} Finite buffer and impermeant pools enter algebraic pH, cell osmolarity, or the neutral charge constraint.\par

\textbf{Constraint.} effective finite buffer pool; must be WT-calibrated independently\par

\textbf{Limit of justification.} An accepted physiological point does not identify this setting or supply a numerical uncertainty interval. No knockout fit or new measurement is performed in this audit.\par

\textbf{Logical domain.} Positive for the retained law; no finite upper physiological bound established.\par

\textbf{Uncertainty status.} genuinely\_unknown\_for\_physiological\_transfer\par

\textbf{Dependence of claims.} Numerical state, current and secretion claims are conditional on this active setting; exact stoichiometric identities are independent of its numerical value. Local numerical dependence is quantified in the linked crosswalk; it establishes neither a measured uncertainty range nor global robustness.\par

\textbf{Discriminating measurement.} Measure intracellular titration capacity together with TIC and cell volume across the experimental pH interval.\par

\textbf{Source.} \nolinkurl{src/modern_full_model/model.py:374}. No independently verified primary measurement attached to this numerical setting; model lineage is the source.\par}

\subsubsection*{Record 19: lumen buffer total fmol}
{\small \nolinkurl{geometry.lumen_buffer_total_fmol}\par

Value: \texttt{1e-09}; units: fmol sites; active: Yes; classification: numerical\_regularisation.\par

Default: 1e-09; differs from default: False.\par

\textbf{Equation role.} Finite buffer and impermeant pools enter algebraic pH, cell osmolarity, or the neutral charge constraint.\par

\textbf{Constraint.} Positive approximation to no fixed luminal buffer. Its size is numerical, not measured.\par

\textbf{Limit of justification.} An accepted physiological point does not identify this setting or supply a numerical uncertainty interval. No knockout fit or new measurement is performed in this audit.\par

\textbf{Logical domain.} Positive for the retained law; no finite upper physiological bound established.\par

\textbf{Uncertainty status.} not\_a\_measured\_uncertainty\_parameter\par

\textbf{Dependence of claims.} Numerical state, current and secretion claims are conditional on this active setting; exact stoichiometric identities are independent of its numerical value.\par

\textbf{Discriminating measurement.} Verify numerical convergence and successful closure on the stated domain; this is not an abundance measurement.\par

\textbf{Source.} \nolinkurl{src/modern_full_model/model.py:374}. No independently verified primary measurement attached to this numerical setting; model lineage is the source.\par}

\subsubsection*{Record 20: fixed cell anion equivalents fmol}
{\small \nolinkurl{geometry.fixed_cell_anion_equivalents_fmol}\par

Value: \texttt{78.06176220594485}; units: fmol charge equivalents; active: Yes; classification: reference\_constraint\_derived.\par

Default: 78.06176220594485; differs from default: False.\par

\textbf{Equation role.} Finite buffer and impermeant pools enter algebraic pH, cell osmolarity, or the neutral charge constraint.\par

\textbf{Constraint.} Derived from the reference initial cations, chloride and total alkalinity. Conditional on that reference tuple, not independently measured.\par

\textbf{Limit of justification.} An accepted physiological point does not identify this setting or supply a numerical uncertainty interval. No knockout fit or new measurement is performed in this audit.\par

\textbf{Logical domain.} Positive for the retained law; no finite upper physiological bound established.\par

\textbf{Uncertainty status.} conditional\_derivation\_no\_experimental\_interval\par

\textbf{Dependence of claims.} Numerical state, current and secretion claims are conditional on this active setting; exact stoichiometric identities are independent of its numerical value.\par

\textbf{Discriminating measurement.} Measure this quantity or its identifiable combination in the same preparation and stimulus protocol, with independent WT data; determine physiological ranges before making global robustness claims.\par

\textbf{Source.} \nolinkurl{src/modern_full_model/model.py:374}. No independently verified primary measurement attached to this numerical setting; model lineage is the source.\par}

\subsubsection*{Record 21: cell impermeant osmoles fmol}
{\small \nolinkurl{geometry.cell_impermeant_osmoles_fmol}\par

Value: \texttt{124.6993693712438}; units: fmol osmoles; active: Yes; classification: inherited\_effective\_setting.\par

Default: 56.16159280366088; differs from default: True.\par

\textbf{Equation role.} Finite buffer and impermeant pools enter algebraic pH, cell osmolarity, or the neutral charge constraint.\par

\textbf{Constraint.} Active R09 background contains 124.69936937124379 fmol, not the dataclass default. An effective osmotic closure setting, not a directly measured molecular pool.\par

\textbf{Limit of justification.} An accepted physiological point does not identify this setting or supply a numerical uncertainty interval. No knockout fit or new measurement is performed in this audit.\par

\textbf{Logical domain.} Positive for the retained law; no finite upper physiological bound established.\par

\textbf{Uncertainty status.} genuinely\_unknown\_for\_physiological\_transfer\par

\textbf{Dependence of claims.} Numerical state, current and secretion claims are conditional on this active setting; exact stoichiometric identities are independent of its numerical value.\par

\textbf{Discriminating measurement.} Measure resting and stimulated cell volume and osmolarity to separate fixed buffer sites from other impermeant osmotic particles.\par

\textbf{Source.} \nolinkurl{src/modern_full_model/model.py:374}. No independently verified primary measurement attached to this numerical setting; model lineage is the source.\par}

\subsubsection*{Record 22: nkcc1 capacity fmol s}
{\small \nolinkurl{homeostasis.nkcc1_capacity_fmol_s}\par

Value: \texttt{0.32}; units: fmol s\textasciicircum{}-1; active: No; classification: inactive\_legacy.\par

Default: 0.08; differs from default: True.\par

\textbf{Equation role.} NHE1 supplies Na and TA; AE2 imports Cl and exports carbon/TA; CO2 transfers carbon only.\par

\textbf{Constraint.} Not used by selected Palk NKCC1 or Cha NHE1 flux. Retained for older models and payload compatibility.\par

\textbf{Limit of justification.} An accepted physiological point does not identify this setting or supply a numerical uncertainty interval. No knockout fit or new measurement is performed in this audit.\par

\textbf{Logical domain.} Positive for the retained law; no finite upper physiological bound established.\par

\textbf{Uncertainty status.} not\_a\_measured\_uncertainty\_parameter\par

\textbf{Dependence of claims.} No independent parent RHS sensitivity. Used only for the explicitly stated metadata, initial reference, design lineage or inverse family role.\par

\textbf{Discriminating measurement.} No new production parameter experiment is implied for this inactive field; measure the actual corresponding state or active pathway instead.\par

\textbf{Source.} \nolinkurl{src/modern_full_model/membranes.py:206}. No independently verified primary measurement attached to this numerical setting; model lineage is the source.\par}

\subsubsection*{Record 23: nhe1 capacity fmol s}
{\small \nolinkurl{homeostasis.nhe1_capacity_fmol_s}\par

Value: \texttt{0.007440980590414482}; units: fmol s\textasciicircum{}-1; active: No; classification: inactive\_legacy.\par

Default: 0.02; differs from default: True.\par

\textbf{Equation role.} NHE1 supplies Na and TA; AE2 imports Cl and exports carbon/TA; CO2 transfers carbon only.\par

\textbf{Constraint.} Not used by selected Palk NKCC1 or Cha NHE1 flux. Retained for older models and payload compatibility.\par

\textbf{Limit of justification.} An accepted physiological point does not identify this setting or supply a numerical uncertainty interval. No knockout fit or new measurement is performed in this audit.\par

\textbf{Logical domain.} Positive for the retained law; no finite upper physiological bound established.\par

\textbf{Uncertainty status.} not\_a\_measured\_uncertainty\_parameter\par

\textbf{Dependence of claims.} No independent parent RHS sensitivity. Used only for the explicitly stated metadata, initial reference, design lineage or inverse family role.\par

\textbf{Discriminating measurement.} No new production parameter experiment is implied for this inactive field; measure the actual corresponding state or active pathway instead.\par

\textbf{Source.} \nolinkurl{src/modern_full_model/membranes.py:206}. No independently verified primary measurement attached to this numerical setting; model lineage is the source.\par}

\subsubsection*{Record 24: nhe1 model}
{\small \nolinkurl{homeostasis.nhe1_model}\par

Value: cha\_2009\_eight\_state\_mod2; units: model identifier; active: Yes; classification: law\_selection.\par

Default: legacy\_tanh; differs from default: True.\par

\textbf{Equation role.} NHE1 supplies Na and TA; AE2 imports Cl and exports carbon/TA; CO2 transfers carbon only.\par

\textbf{Constraint.} Production explicitly selects Cha 2009 eight-state Mod2; the default field note refers to an older selection.\par

\textbf{Limit of justification.} An accepted physiological point does not identify this setting or supply a numerical uncertainty interval. No knockout fit or new measurement is performed in this audit.\par

\textbf{Logical domain.} Selected discrete setting or stated numerical reference.\par

\textbf{Uncertainty status.} not\_a\_measured\_uncertainty\_parameter\par

\textbf{Dependence of claims.} Numerical state, current and secretion claims are conditional on this active setting; exact stoichiometric identities are independent of its numerical value.\par

\textbf{Discriminating measurement.} Measure this quantity or its identifiable combination in the same preparation and stimulus protocol, with independent WT data; determine physiological ranges before making global robustness claims.\par

\textbf{Source.} \nolinkurl{src/modern_full_model/membranes.py:206}. No independently verified primary measurement attached to this numerical setting; model lineage is the source.\par}

\subsubsection*{Record 25: nhe1 published g fmol s}
{\small \nolinkurl{homeostasis.nhe1_published_g_fmol_s}\par

Value: \texttt{0.0305}; units: fmol s\textasciicircum{}-1; active: No; classification: inactive\_legacy.\par

Default: 0.0305; differs from default: False.\par

\textbf{Equation role.} NHE1 supplies Na and TA; AE2 imports Cl and exports carbon/TA; CO2 transfers carbon only.\par

\textbf{Constraint.} Not used by selected Palk NKCC1 or Cha NHE1 flux. Retained for older models and payload compatibility.\par

\textbf{Limit of justification.} An accepted physiological point does not identify this setting or supply a numerical uncertainty interval. No knockout fit or new measurement is performed in this audit.\par

\textbf{Logical domain.} Positive for the retained law; no finite upper physiological bound established.\par

\textbf{Uncertainty status.} not\_a\_measured\_uncertainty\_parameter\par

\textbf{Dependence of claims.} No independent parent RHS sensitivity. Used only for the explicitly stated metadata, initial reference, design lineage or inverse family role.\par

\textbf{Discriminating measurement.} No new production parameter experiment is implied for this inactive field; measure the actual corresponding state or active pathway instead.\par

\textbf{Source.} \nolinkurl{src/modern_full_model/membranes.py:206}. No independently verified primary measurement attached to this numerical setting; model lineage is the source.\par}

\subsubsection*{Record 26: nhe1 published k h mM}
{\small \nolinkurl{homeostasis.nhe1_published_k_h_mM}\par

Value: \texttt{0.00045}; units: mM; active: No; classification: inactive\_legacy.\par

Default: 0.00045; differs from default: False.\par

\textbf{Equation role.} NHE1 supplies Na and TA; AE2 imports Cl and exports carbon/TA; CO2 transfers carbon only.\par

\textbf{Constraint.} Not used by selected Palk NKCC1 or Cha NHE1 flux. Retained for older models and payload compatibility.\par

\textbf{Limit of justification.} An accepted physiological point does not identify this setting or supply a numerical uncertainty interval. No knockout fit or new measurement is performed in this audit.\par

\textbf{Logical domain.} Positive for the retained law; no finite upper physiological bound established.\par

\textbf{Uncertainty status.} not\_a\_measured\_uncertainty\_parameter\par

\textbf{Dependence of claims.} No independent parent RHS sensitivity. Used only for the explicitly stated metadata, initial reference, design lineage or inverse family role.\par

\textbf{Discriminating measurement.} No new production parameter experiment is implied for this inactive field; measure the actual corresponding state or active pathway instead.\par

\textbf{Source.} \nolinkurl{src/modern_full_model/membranes.py:206}. No independently verified primary measurement attached to this numerical setting; model lineage is the source.\par}

\subsubsection*{Record 27: nhe1 published k na mM}
{\small \nolinkurl{homeostasis.nhe1_published_k_na_mM}\par

Value: \texttt{15}; units: mM; active: No; classification: inactive\_legacy.\par

Default: 15.0; differs from default: False.\par

\textbf{Equation role.} NHE1 supplies Na and TA; AE2 imports Cl and exports carbon/TA; CO2 transfers carbon only.\par

\textbf{Constraint.} Not used by selected Palk NKCC1 or Cha NHE1 flux. Retained for older models and payload compatibility.\par

\textbf{Limit of justification.} An accepted physiological point does not identify this setting or supply a numerical uncertainty interval. No knockout fit or new measurement is performed in this audit.\par

\textbf{Logical domain.} Positive for the retained law; no finite upper physiological bound established.\par

\textbf{Uncertainty status.} not\_a\_measured\_uncertainty\_parameter\par

\textbf{Dependence of claims.} No independent parent RHS sensitivity. Used only for the explicitly stated metadata, initial reference, design lineage or inverse family role.\par

\textbf{Discriminating measurement.} No new production parameter experiment is implied for this inactive field; measure the actual corresponding state or active pathway instead.\par

\textbf{Source.} \nolinkurl{src/modern_full_model/membranes.py:206}. No independently verified primary measurement attached to this numerical setting; model lineage is the source.\par}

\subsubsection*{Record 28: nhe1 cha carrier amount fmol}
{\small \nolinkurl{homeostasis.nhe1_cha_carrier_amount_fmol}\par

Value: \texttt{2.339370005697548e-05}; units: fmol of NHE1 carriers; active: Yes; classification: WT\_constraint\_derived.\par

Default: None; differs from default: True.\par

\textbf{Equation role.} NHE1 supplies Na and TA; AE2 imports Cl and exports carbon/TA; CO2 transfers carbon only.\par

\textbf{Constraint.} One salivary carrier amount calibrated against a resting pH target 6.91 on the selected R09 background. It is not the cardiac carrier count or a direct salivary abundance measurement.\par

\textbf{Limit of justification.} An accepted physiological point does not identify this setting or supply a numerical uncertainty interval. No knockout fit or new measurement is performed in this audit.\par

\textbf{Logical domain.} Positive for the retained law; no finite upper physiological bound established.\par

\textbf{Uncertainty status.} conditional\_derivation\_no\_experimental\_interval\par

\textbf{Dependence of claims.} Numerical state, current and secretion claims are conditional on this active setting; exact stoichiometric identities are independent of its numerical value. Local numerical dependence is quantified in the linked crosswalk; it establishes neither a measured uncertainty range nor global robustness.\par

\textbf{Discriminating measurement.} Measure resting and stimulated Na/H exchange turnover during controlled acid loads, with pH and cell volume; distinguish carrier abundance from kinetic activity.\par

\textbf{Source.} \nolinkurl{src/modern_full_model/membranes.py:206}. No independently verified primary measurement attached to this numerical setting; model lineage is the source.\par}

\subsubsection*{Record 29: ae2 capacity fmol s}
{\small \nolinkurl{homeostasis.ae2_capacity_fmol_s}\par

Value: \texttt{0.005}; units: fmol s\textasciicircum{}-1; active: Yes; classification: effective\_assumption.\par

Default: 0.005; differs from default: False.\par

\textbf{Equation role.} NHE1 supplies Na and TA; AE2 imports Cl and exports carbon/TA; CO2 transfers carbon only.\par

\textbf{Constraint.} small placeholder capacity; AE2-null is a validation gate\par

\textbf{Limit of justification.} An accepted physiological point does not identify this setting or supply a numerical uncertainty interval. No knockout fit or new measurement is performed in this audit.\par

\textbf{Logical domain.} Positive for the retained law; no finite upper physiological bound established.\par

\textbf{Uncertainty status.} genuinely\_unknown\_for\_physiological\_transfer\par

\textbf{Dependence of claims.} Numerical state, current and secretion claims are conditional on this active setting; exact stoichiometric identities are independent of its numerical value. Local numerical dependence is quantified in the linked crosswalk; it establishes neither a measured uncertainty range nor global robustness.\par

\textbf{Discriminating measurement.} Measure this quantity or its identifiable combination in the same preparation and stimulus protocol, with independent WT data; determine physiological ranges before making global robustness claims.\par

\textbf{Source.} \nolinkurl{src/modern_full_model/membranes.py:206}. No independently verified primary measurement attached to this numerical setting; model lineage is the source.\par}

\subsubsection*{Record 30: thermodynamic saturation log width}
{\small \nolinkurl{homeostasis.thermodynamic_saturation_log_width}\par

Value: \texttt{2}; units: dimensionless log-affinity; active: Yes; classification: effective\_assumption.\par

Default: 2.0; differs from default: False.\par

\textbf{Equation role.} NHE1 supplies Na and TA; AE2 imports Cl and exports carbon/TA; CO2 transfers carbon only.\par

\textbf{Constraint.} sets tanh approach to capacity without moving reversal\par

\textbf{Limit of justification.} An accepted physiological point does not identify this setting or supply a numerical uncertainty interval. No knockout fit or new measurement is performed in this audit.\par

\textbf{Logical domain.} Positive for the retained law; no finite upper physiological bound established.\par

\textbf{Uncertainty status.} genuinely\_unknown\_for\_physiological\_transfer\par

\textbf{Dependence of claims.} Numerical state, current and secretion claims are conditional on this active setting; exact stoichiometric identities are independent of its numerical value.\par

\textbf{Discriminating measurement.} Measure this quantity or its identifiable combination in the same preparation and stimulus protocol, with independent WT data; determine physiological ranges before making global robustness claims.\par

\textbf{Source.} \nolinkurl{src/modern_full_model/membranes.py:206}. No independently verified primary measurement attached to this numerical setting; model lineage is the source.\par}

\subsubsection*{Record 31: co2 basolateral permeability fmol s mM}
{\small \nolinkurl{homeostasis.co2_basolateral_permeability_fmol_s_mM}\par

Value: \texttt{0.02}; units: fmol s\textasciicircum{}-1 mM\textasciicircum{}-1; active: Yes; classification: effective\_assumption.\par

Default: 0.02; differs from default: False.\par

\textbf{Equation role.} J\_CO2,b=P\_b([CO2]bath-[CO2]cell); J\_CO2,a=P\_a([CO2]lumen-[CO2]cell); modifies carbon amounts, not alkalinity.\par

\textbf{Constraint.} linear neutral CO2 exchange coefficient\par

\textbf{Limit of justification.} An accepted physiological point does not identify this setting or supply a numerical uncertainty interval. No knockout fit or new measurement is performed in this audit.\par

\textbf{Logical domain.} Positive for the retained law; no finite upper physiological bound established.\par

\textbf{Uncertainty status.} genuinely\_unknown\_for\_physiological\_transfer\par

\textbf{Dependence of claims.} Numerical state, current and secretion claims are conditional on this active setting; exact stoichiometric identities are independent of its numerical value. Local numerical dependence is quantified in the linked crosswalk; it establishes neither a measured uncertainty range nor global robustness.\par

\textbf{Discriminating measurement.} Measure this quantity or its identifiable combination in the same preparation and stimulus protocol, with independent WT data; determine physiological ranges before making global robustness claims.\par

\textbf{Source.} \nolinkurl{src/modern_full_model/nbc_minimal.py:589}. No independently verified primary measurement attached to this numerical setting; model lineage is the source.\par}

\subsubsection*{Record 32: co2 apical permeability fmol s mM}
{\small \nolinkurl{homeostasis.co2_apical_permeability_fmol_s_mM}\par

Value: \texttt{0.01}; units: fmol s\textasciicircum{}-1 mM\textasciicircum{}-1; active: Yes; classification: effective\_assumption.\par

Default: 0.01; differs from default: False.\par

\textbf{Equation role.} J\_CO2,b=P\_b([CO2]bath-[CO2]cell); J\_CO2,a=P\_a([CO2]lumen-[CO2]cell); modifies carbon amounts, not alkalinity.\par

\textbf{Constraint.} linear neutral CO2 exchange coefficient\par

\textbf{Limit of justification.} An accepted physiological point does not identify this setting or supply a numerical uncertainty interval. No knockout fit or new measurement is performed in this audit.\par

\textbf{Logical domain.} Positive for the retained law; no finite upper physiological bound established.\par

\textbf{Uncertainty status.} genuinely\_unknown\_for\_physiological\_transfer\par

\textbf{Dependence of claims.} Numerical state, current and secretion claims are conditional on this active setting; exact stoichiometric identities are independent of its numerical value.\par

\textbf{Discriminating measurement.} Measure this quantity or its identifiable combination in the same preparation and stimulus protocol, with independent WT data; determine physiological ranges before making global robustness claims.\par

\textbf{Source.} \nolinkurl{src/modern_full_model/nbc_minimal.py:589}. No independently verified primary measurement attached to this numerical setting; model lineage is the source.\par}

\subsubsection*{Record 33: nak capacity fmol s}
{\small \nolinkurl{membranes.nak_capacity_fmol_s}\par

Value: \texttt{0.08}; units: fmol cycles s\textasciicircum{}-1; active: Yes; classification: effective\_assumption.\par

Default: 0.08; differs from default: False.\par

\textbf{Equation role.} Pump and channel currents determine two algebraic voltages and conserved cell/lumen source vectors.\par

\textbf{Constraint.} total pump capacity shared between membranes\par

\textbf{Limit of justification.} An accepted physiological point does not identify this setting or supply a numerical uncertainty interval. No knockout fit or new measurement is performed in this audit.\par

\textbf{Logical domain.} Positive for the retained law; no finite upper physiological bound established.\par

\textbf{Uncertainty status.} genuinely\_unknown\_for\_physiological\_transfer\par

\textbf{Dependence of claims.} Numerical state, current and secretion claims are conditional on this active setting; exact stoichiometric identities are independent of its numerical value. Local numerical dependence is quantified in the linked crosswalk; it establishes neither a measured uncertainty range nor global robustness.\par

\textbf{Discriminating measurement.} Measure pump current or turnover versus intracellular Na during the secretion protocol and its apical distribution.\par

\textbf{Source.} \nolinkurl{src/modern_full_model/nbc_minimal.py:248}. No independently verified primary measurement attached to this numerical setting; model lineage is the source.\par}

\subsubsection*{Record 34: nak na half mM}
{\small \nolinkurl{membranes.nak_na_half_mM}\par

Value: \texttt{10}; units: mM; active: Yes; classification: inherited\_model\_setting.\par

Default: 10.0; differs from default: False.\par

\textbf{Equation role.} Pump and channel currents determine two algebraic voltages and conserved cell/lumen source vectors.\par

\textbf{Constraint.} effective intracellular Na half-saturation starting value Source metadata identifies a model lineage; independent experimental measurement is not established here.\par

\textbf{Limit of justification.} An accepted physiological point does not identify this setting or supply a numerical uncertainty interval. No knockout fit or new measurement is performed in this audit.\par

\textbf{Logical domain.} Positive for the retained law; no finite upper physiological bound established.\par

\textbf{Uncertainty status.} genuinely\_unknown\_for\_physiological\_transfer\par

\textbf{Dependence of claims.} Numerical state, current and secretion claims are conditional on this active setting; exact stoichiometric identities are independent of its numerical value.\par

\textbf{Discriminating measurement.} Measure this quantity or its identifiable combination in the same preparation and stimulus protocol, with independent WT data; determine physiological ranges before making global robustness claims.\par

\textbf{Source.} \nolinkurl{src/modern_full_model/nbc_minimal.py:248}. No independently verified primary measurement attached to this numerical setting; model lineage is the source.\par}

\subsubsection*{Record 35: nak k half mM}
{\small \nolinkurl{membranes.nak_k_half_mM}\par

Value: \texttt{1.5}; units: mM; active: Yes; classification: inherited\_model\_setting.\par

Default: 1.5; differs from default: False.\par

\textbf{Equation role.} Pump and channel currents determine two algebraic voltages and conserved cell/lumen source vectors.\par

\textbf{Constraint.} effective extracellular K half-saturation starting value Source metadata identifies a model lineage; independent experimental measurement is not established here.\par

\textbf{Limit of justification.} An accepted physiological point does not identify this setting or supply a numerical uncertainty interval. No knockout fit or new measurement is performed in this audit.\par

\textbf{Logical domain.} Positive for the retained law; no finite upper physiological bound established.\par

\textbf{Uncertainty status.} genuinely\_unknown\_for\_physiological\_transfer\par

\textbf{Dependence of claims.} Numerical state, current and secretion claims are conditional on this active setting; exact stoichiometric identities are independent of its numerical value.\par

\textbf{Discriminating measurement.} Measure this quantity or its identifiable combination in the same preparation and stimulus protocol, with independent WT data; determine physiological ranges before making global robustness claims.\par

\textbf{Source.} \nolinkurl{src/modern_full_model/nbc_minimal.py:248}. No independently verified primary measurement attached to this numerical setting; model lineage is the source.\par}

\subsubsection*{Record 36: apical pump fraction}
{\small \nolinkurl{membranes.apical_pump_fraction}\par

Value: \texttt{0.075075}; units: fraction; active: Yes; classification: effective\_assumption.\par

Default: 0.1; differs from default: True.\par

\textbf{Equation role.} Pump and channel currents determine two algebraic voltages and conserved cell/lumen source vectors.\par

\textbf{Constraint.} apical localization is supported; quantitative fraction is not measured\par

\textbf{Limit of justification.} An accepted physiological point does not identify this setting or supply a numerical uncertainty interval. No knockout fit or new measurement is performed in this audit.\par

\textbf{Logical domain.} [0,1] by fractional definition only; not a measured uncertainty interval.\par

\textbf{Uncertainty status.} genuinely\_unknown\_for\_physiological\_transfer\par

\textbf{Dependence of claims.} Numerical state, current and secretion claims are conditional on this active setting; exact stoichiometric identities are independent of its numerical value. Local numerical dependence is quantified in the linked crosswalk; it establishes neither a measured uncertainty range nor global robustness.\par

\textbf{Discriminating measurement.} Quantify apical versus basolateral functional pump turnover, rather than inferring the fraction from localisation alone.\par

\textbf{Source.} \nolinkurl{src/modern_full_model/nbc_minimal.py:248}. No independently verified primary measurement attached to this numerical setting; model lineage is the source.\par}

\subsubsection*{Record 37: apical k fraction}
{\small \nolinkurl{membranes.apical_k_fraction}\par

Value: \texttt{0.3}; units: fraction; active: Yes; classification: effective\_assumption.\par

Default: 0.1; differs from default: True.\par

\textbf{Equation role.} Pump and channel currents determine two algebraic voltages and conserved cell/lumen source vectors.\par

\textbf{Constraint.} apical K current is supported; quantitative fraction is not measured\par

\textbf{Limit of justification.} An accepted physiological point does not identify this setting or supply a numerical uncertainty interval. No knockout fit or new measurement is performed in this audit.\par

\textbf{Logical domain.} [0,1] by fractional definition only; not a measured uncertainty interval.\par

\textbf{Uncertainty status.} genuinely\_unknown\_for\_physiological\_transfer\par

\textbf{Dependence of claims.} Numerical state, current and secretion claims are conditional on this active setting; exact stoichiometric identities are independent of its numerical value.\par

\textbf{Discriminating measurement.} Measure this quantity or its identifiable combination in the same preparation and stimulus protocol, with independent WT data; determine physiological ranges before making global robustness claims.\par

\textbf{Source.} \nolinkurl{src/modern_full_model/nbc_minimal.py:248}. No independently verified primary measurement attached to this numerical setting; model lineage is the source.\par}

\subsubsection*{Record 38: g k total S}
{\small \nolinkurl{membranes.g_k_total_S}\par

Value: \texttt{1.4e-08}; units: S; active: Yes; classification: effective\_assumption.\par

Default: 4e-10; differs from default: True.\par

\textbf{Equation role.} Pump and channel currents determine two algebraic voltages and conserved cell/lumen source vectors.\par

\textbf{Constraint.} total Ca-activated K conductance shared between membranes\par

\textbf{Limit of justification.} An accepted physiological point does not identify this setting or supply a numerical uncertainty interval. No knockout fit or new measurement is performed in this audit.\par

\textbf{Logical domain.} Positive for the retained law; no finite upper physiological bound established.\par

\textbf{Uncertainty status.} genuinely\_unknown\_for\_physiological\_transfer\par

\textbf{Dependence of claims.} Numerical state, current and secretion claims are conditional on this active setting; exact stoichiometric identities are independent of its numerical value.\par

\textbf{Discriminating measurement.} Measure this quantity or its identifiable combination in the same preparation and stimulus protocol, with independent WT data; determine physiological ranges before making global robustness claims.\par

\textbf{Source.} \nolinkurl{src/modern_full_model/nbc_minimal.py:248}. No independently verified primary measurement attached to this numerical setting; model lineage is the source.\par}

\subsubsection*{Record 39: g cl apical S}
{\small \nolinkurl{membranes.g_cl_apical_S}\par

Value: \texttt{3.14e-08}; units: S; active: Yes; classification: effective\_assumption.\par

Default: 6e-10; differs from default: True.\par

\textbf{Equation role.} Pump and channel currents determine two algebraic voltages and conserved cell/lumen source vectors.\par

\textbf{Constraint.} apical Ca-activated Cl conductance\par

\textbf{Limit of justification.} An accepted physiological point does not identify this setting or supply a numerical uncertainty interval. No knockout fit or new measurement is performed in this audit.\par

\textbf{Logical domain.} Positive for the retained law; no finite upper physiological bound established.\par

\textbf{Uncertainty status.} genuinely\_unknown\_for\_physiological\_transfer\par

\textbf{Dependence of claims.} Numerical state, current and secretion claims are conditional on this active setting; exact stoichiometric identities are independent of its numerical value. Local numerical dependence is quantified in the linked crosswalk; it establishes neither a measured uncertainty range nor global robustness.\par

\textbf{Discriminating measurement.} Measure CaCC current versus calcium and voltage in matched WT and AE4 deficient cells, including onset and sustained recruitment.\par

\textbf{Source.} \nolinkurl{src/modern_full_model/nbc_minimal.py:248}. No independently verified primary measurement attached to this numerical setting; model lineage is the source.\par}

\subsubsection*{Record 40: calcium half uM}
{\small \nolinkurl{membranes.calcium_half_uM}\par

Value: \texttt{0.26}; units: uM; active: Yes; classification: effective\_assumption.\par

Default: 0.3; differs from default: True.\par

\textbf{Equation role.} Pump and channel currents determine two algebraic voltages and conserved cell/lumen source vectors.\par

\textbf{Constraint.} common test gating scale, not a frozen WT estimate\par

\textbf{Limit of justification.} An accepted physiological point does not identify this setting or supply a numerical uncertainty interval. No knockout fit or new measurement is performed in this audit.\par

\textbf{Logical domain.} Positive for the retained law; no finite upper physiological bound established.\par

\textbf{Uncertainty status.} genuinely\_unknown\_for\_physiological\_transfer\par

\textbf{Dependence of claims.} Numerical state, current and secretion claims are conditional on this active setting; exact stoichiometric identities are independent of its numerical value.\par

\textbf{Discriminating measurement.} Measure this quantity or its identifiable combination in the same preparation and stimulus protocol, with independent WT data; determine physiological ranges before making global robustness claims.\par

\textbf{Source.} \nolinkurl{src/modern_full_model/nbc_minimal.py:248}. No independently verified primary measurement attached to this numerical setting; model lineage is the source.\par}

\subsubsection*{Record 41: calcium hill}
{\small \nolinkurl{membranes.calcium_hill}\par

Value: \texttt{1.46}; units: dimensionless; active: Yes; classification: effective\_assumption.\par

Default: 2.0; differs from default: True.\par

\textbf{Equation role.} Pump and channel currents determine two algebraic voltages and conserved cell/lumen source vectors.\par

\textbf{Constraint.} effective Ca gating exponent\par

\textbf{Limit of justification.} An accepted physiological point does not identify this setting or supply a numerical uncertainty interval. No knockout fit or new measurement is performed in this audit.\par

\textbf{Logical domain.} Positive for the retained law; no finite upper physiological bound established.\par

\textbf{Uncertainty status.} genuinely\_unknown\_for\_physiological\_transfer\par

\textbf{Dependence of claims.} Numerical state, current and secretion claims are conditional on this active setting; exact stoichiometric identities are independent of its numerical value.\par

\textbf{Discriminating measurement.} Measure this quantity or its identifiable combination in the same preparation and stimulus protocol, with independent WT data; determine physiological ranges before making global robustness claims.\par

\textbf{Source.} \nolinkurl{src/modern_full_model/nbc_minimal.py:248}. No independently verified primary measurement attached to this numerical setting; model lineage is the source.\par}

\subsubsection*{Record 42: g para na S}
{\small \nolinkurl{membranes.g_para_na_S}\par

Value: \texttt{3e-10}; units: S; active: Yes; classification: inherited\_model\_setting.\par

Default: 3e-10; differs from default: False.\par

\textbf{Equation role.} Pump and channel currents determine two algebraic voltages and conserved cell/lumen source vectors.\par

\textbf{Constraint.} historical-lineage paracellular Na pathway; value not frozen Source metadata identifies a model lineage; independent experimental measurement is not established here.\par

\textbf{Limit of justification.} An accepted physiological point does not identify this setting or supply a numerical uncertainty interval. No knockout fit or new measurement is performed in this audit.\par

\textbf{Logical domain.} Positive for the retained law; no finite upper physiological bound established.\par

\textbf{Uncertainty status.} genuinely\_unknown\_for\_physiological\_transfer\par

\textbf{Dependence of claims.} Numerical state, current and secretion claims are conditional on this active setting; exact stoichiometric identities are independent of its numerical value.\par

\textbf{Discriminating measurement.} Measure this quantity or its identifiable combination in the same preparation and stimulus protocol, with independent WT data; determine physiological ranges before making global robustness claims.\par

\textbf{Source.} \nolinkurl{src/modern_full_model/nbc_minimal.py:248}. No independently verified primary measurement attached to this numerical setting; model lineage is the source.\par}

\subsubsection*{Record 43: g para k S}
{\small \nolinkurl{membranes.g_para_k_S}\par

Value: \texttt{5e-11}; units: S; active: Yes; classification: inherited\_model\_setting.\par

Default: 5e-11; differs from default: False.\par

\textbf{Equation role.} Pump and channel currents determine two algebraic voltages and conserved cell/lumen source vectors.\par

\textbf{Constraint.} historical-lineage paracellular K pathway; value not frozen Source metadata identifies a model lineage; independent experimental measurement is not established here.\par

\textbf{Limit of justification.} An accepted physiological point does not identify this setting or supply a numerical uncertainty interval. No knockout fit or new measurement is performed in this audit.\par

\textbf{Logical domain.} Positive for the retained law; no finite upper physiological bound established.\par

\textbf{Uncertainty status.} genuinely\_unknown\_for\_physiological\_transfer\par

\textbf{Dependence of claims.} Numerical state, current and secretion claims are conditional on this active setting; exact stoichiometric identities are independent of its numerical value.\par

\textbf{Discriminating measurement.} Measure this quantity or its identifiable combination in the same preparation and stimulus protocol, with independent WT data; determine physiological ranges before making global robustness claims.\par

\textbf{Source.} \nolinkurl{src/modern_full_model/nbc_minimal.py:248}. No independently verified primary measurement attached to this numerical setting; model lineage is the source.\par}

\subsubsection*{Record 44: g para cl S}
{\small \nolinkurl{membranes.g_para_cl_S}\par

Value: \texttt{2.5e-10}; units: S; active: Yes; classification: effective\_assumption.\par

Default: 2.5e-10; differs from default: False.\par

\textbf{Equation role.} Pump and channel currents determine two algebraic voltages and conserved cell/lumen source vectors.\par

\textbf{Constraint.} charge-complete paracellular Cl pathway\par

\textbf{Limit of justification.} An accepted physiological point does not identify this setting or supply a numerical uncertainty interval. No knockout fit or new measurement is performed in this audit.\par

\textbf{Logical domain.} Positive for the retained law; no finite upper physiological bound established.\par

\textbf{Uncertainty status.} genuinely\_unknown\_for\_physiological\_transfer\par

\textbf{Dependence of claims.} Numerical state, current and secretion claims are conditional on this active setting; exact stoichiometric identities are independent of its numerical value.\par

\textbf{Discriminating measurement.} Measure this quantity or its identifiable combination in the same preparation and stimulus protocol, with independent WT data; determine physiological ranges before making global robustness claims.\par

\textbf{Source.} \nolinkurl{src/modern_full_model/nbc_minimal.py:248}. No independently verified primary measurement attached to this numerical setting; model lineage is the source.\par}

\subsubsection*{Record 45: g para hco3 S}
{\small \nolinkurl{membranes.g_para_hco3_S}\par

Value: \texttt{5e-11}; units: S; active: Yes; classification: effective\_assumption.\par

Default: 5e-11; differs from default: False.\par

\textbf{Equation role.} Pump and channel currents determine two algebraic voltages and conserved cell/lumen source vectors.\par

\textbf{Constraint.} charge-complete paracellular bicarbonate pathway\par

\textbf{Limit of justification.} An accepted physiological point does not identify this setting or supply a numerical uncertainty interval. No knockout fit or new measurement is performed in this audit.\par

\textbf{Logical domain.} Positive for the retained law; no finite upper physiological bound established.\par

\textbf{Uncertainty status.} genuinely\_unknown\_for\_physiological\_transfer\par

\textbf{Dependence of claims.} Numerical state, current and secretion claims are conditional on this active setting; exact stoichiometric identities are independent of its numerical value.\par

\textbf{Discriminating measurement.} Measure this quantity or its identifiable combination in the same preparation and stimulus protocol, with independent WT data; determine physiological ranges before making global robustness claims.\par

\textbf{Source.} \nolinkurl{src/modern_full_model/nbc_minimal.py:248}. No independently verified primary measurement attached to this numerical setting; model lineage is the source.\par}

\subsubsection*{Record 46: g apical background S}
{\small \nolinkurl{membranes.g_apical_background_S}\par

Value: \texttt{0}; units: S; active: Yes; classification: effective\_assumption.\par

Default: 0.0; differs from default: False.\par

\textbf{Equation role.} Pump and channel currents determine two algebraic voltages and conserved cell/lumen source vectors.\par

\textbf{Constraint.} zero preserves the exact no-apical-K nesting limit An exact zero is an explicit structural omission, not a small positive fitted value.\par

\textbf{Limit of justification.} An accepted physiological point does not identify this setting or supply a numerical uncertainty interval. No knockout fit or new measurement is performed in this audit.\par

\textbf{Logical domain.} Exact zero records a declared omitted contribution or protocol normalisation; it is not a confidence bound.\par

\textbf{Uncertainty status.} genuinely\_unknown\_for\_physiological\_transfer\par

\textbf{Dependence of claims.} Numerical state, current and secretion claims are conditional on this active setting; exact stoichiometric identities are independent of its numerical value.\par

\textbf{Discriminating measurement.} Measure this quantity or its identifiable combination in the same preparation and stimulus protocol, with independent WT data; determine physiological ranges before making global robustness claims.\par

\textbf{Source.} \nolinkurl{src/modern_full_model/nbc_minimal.py:248}. No independently verified primary measurement attached to this numerical setting; model lineage is the source.\par}

\subsubsection*{Record 47: g basolateral background S}
{\small \nolinkurl{membranes.g_basolateral_background_S}\par

Value: \texttt{0}; units: S; active: Yes; classification: effective\_assumption.\par

Default: 5e-12; differs from default: True.\par

\textbf{Equation role.} Pump and channel currents determine two algebraic voltages and conserved cell/lumen source vectors.\par

\textbf{Constraint.} small K-selective regularizer; set zero only if closure remains nonsingular An exact zero is an explicit structural omission, not a small positive fitted value.\par

\textbf{Limit of justification.} An accepted physiological point does not identify this setting or supply a numerical uncertainty interval. No knockout fit or new measurement is performed in this audit.\par

\textbf{Logical domain.} Exact zero records a declared omitted contribution or protocol normalisation; it is not a confidence bound.\par

\textbf{Uncertainty status.} genuinely\_unknown\_for\_physiological\_transfer\par

\textbf{Dependence of claims.} Numerical state, current and secretion claims are conditional on this active setting; exact stoichiometric identities are independent of its numerical value.\par

\textbf{Discriminating measurement.} Measure this quantity or its identifiable combination in the same preparation and stimulus protocol, with independent WT data; determine physiological ranges before making global robustness claims.\par

\textbf{Source.} \nolinkurl{src/modern_full_model/nbc_minimal.py:248}. No independently verified primary measurement attached to this numerical setting; model lineage is the source.\par}

\subsubsection*{Record 48: apical hydraulic pL s mOsm}
{\small \nolinkurl{water.apical_hydraulic_pL_s_mOsm}\par

Value: \texttt{0.00432}; units: pL s\textasciicircum{}-1 (mOsm L\textasciicircum{}-1)\textasciicircum{}-1; active: Yes; classification: published\_model\_conversion.\par

Default: 0.00432; differs from default: False.\par

\textbf{Equation role.} Q\_b=L\_b(O\_i-O\_b), Q\_a=L\_a(O\_l-O\_i), Q\_p=L\_p(O\_l-O\_b), Q\_out=k\_out max(V\_l-V\_dead,0).\par

\textbf{Constraint.} Palk 2010/2018 Pa=4.32e-12 L\textasciicircum{}2 mol\textasciicircum{}-1 s\textasciicircum{}-1 converted by 1 mM=1e-3 mol/L and 1 L=1e12 pL Inherited model coefficient, not a new direct measurement.\par

\textbf{Limit of justification.} An accepted physiological point does not identify this setting or supply a numerical uncertainty interval. No knockout fit or new measurement is performed in this audit.\par

\textbf{Logical domain.} Positive for the retained law; no finite upper physiological bound established.\par

\textbf{Uncertainty status.} genuinely\_unknown\_for\_physiological\_transfer\par

\textbf{Dependence of claims.} Numerical state, current and secretion claims are conditional on this active setting; exact stoichiometric identities are independent of its numerical value. Local numerical dependence is quantified in the linked crosswalk; it establishes neither a measured uncertainty range nor global robustness.\par

\textbf{Discriminating measurement.} Measure this quantity or its identifiable combination in the same preparation and stimulus protocol, with independent WT data; determine physiological ranges before making global robustness claims.\par

\textbf{Source.} \nolinkurl{src/modern_full_model/water.py:93}. Palk 2010/2018 inherited whole cell water coefficients; dimensional conversion is recorded in parameters.py.\par}

\subsubsection*{Record 49: basolateral hydraulic pL s mOsm}
{\small \nolinkurl{water.basolateral_hydraulic_pL_s_mOsm}\par

Value: \texttt{0.0515}; units: pL s\textasciicircum{}-1 (mOsm L\textasciicircum{}-1)\textasciicircum{}-1; active: Yes; classification: published\_model\_conversion.\par

Default: 0.0515; differs from default: False.\par

\textbf{Equation role.} Q\_b=L\_b(O\_i-O\_b), Q\_a=L\_a(O\_l-O\_i), Q\_p=L\_p(O\_l-O\_b), Q\_out=k\_out max(V\_l-V\_dead,0).\par

\textbf{Constraint.} Palk 2010/2018 Pb=5.15e-11 L\textasciicircum{}2 mol\textasciicircum{}-1 s\textasciicircum{}-1 converted by 1 mM=1e-3 mol/L and 1 L=1e12 pL Inherited model coefficient, not a new direct measurement.\par

\textbf{Limit of justification.} An accepted physiological point does not identify this setting or supply a numerical uncertainty interval. No knockout fit or new measurement is performed in this audit.\par

\textbf{Logical domain.} Positive for the retained law; no finite upper physiological bound established.\par

\textbf{Uncertainty status.} genuinely\_unknown\_for\_physiological\_transfer\par

\textbf{Dependence of claims.} Numerical state, current and secretion claims are conditional on this active setting; exact stoichiometric identities are independent of its numerical value.\par

\textbf{Discriminating measurement.} Measure this quantity or its identifiable combination in the same preparation and stimulus protocol, with independent WT data; determine physiological ranges before making global robustness claims.\par

\textbf{Source.} \nolinkurl{src/modern_full_model/water.py:93}. Palk 2010/2018 inherited whole cell water coefficients; dimensional conversion is recorded in parameters.py.\par}

\subsubsection*{Record 50: paracellular hydraulic pL s mOsm}
{\small \nolinkurl{water.paracellular_hydraulic_pL_s_mOsm}\par

Value: \texttt{0.00026}; units: pL s\textasciicircum{}-1 (mOsm L\textasciicircum{}-1)\textasciicircum{}-1; active: Yes; classification: published\_model\_conversion.\par

Default: 0.00026; differs from default: False.\par

\textbf{Equation role.} Q\_b=L\_b(O\_i-O\_b), Q\_a=L\_a(O\_l-O\_i), Q\_p=L\_p(O\_l-O\_b), Q\_out=k\_out max(V\_l-V\_dead,0).\par

\textbf{Constraint.} Palk 2010/2018 Pt=2.6e-13 L\textasciicircum{}2 mol\textasciicircum{}-1 s\textasciicircum{}-1 converted by 1 mM=1e-3 mol/L and 1 L=1e12 pL Inherited model coefficient, not a new direct measurement.\par

\textbf{Limit of justification.} An accepted physiological point does not identify this setting or supply a numerical uncertainty interval. No knockout fit or new measurement is performed in this audit.\par

\textbf{Logical domain.} Positive for the retained law; no finite upper physiological bound established.\par

\textbf{Uncertainty status.} genuinely\_unknown\_for\_physiological\_transfer\par

\textbf{Dependence of claims.} Numerical state, current and secretion claims are conditional on this active setting; exact stoichiometric identities are independent of its numerical value.\par

\textbf{Discriminating measurement.} Measure this quantity or its identifiable combination in the same preparation and stimulus protocol, with independent WT data; determine physiological ranges before making global robustness claims.\par

\textbf{Source.} \nolinkurl{src/modern_full_model/water.py:93}. Palk 2010/2018 inherited whole cell water coefficients; dimensional conversion is recorded in parameters.py.\par}

\subsubsection*{Record 51: lumen dead volume pL}
{\small \nolinkurl{water.lumen_dead_volume_pL}\par

Value: \texttt{0.1}; units: pL; active: Yes; classification: effective\_assumption.\par

Default: 0.1; differs from default: False.\par

\textbf{Equation role.} Q\_b=L\_b(O\_i-O\_b), Q\_a=L\_a(O\_l-O\_i), Q\_p=L\_p(O\_l-O\_b), Q\_out=k\_out max(V\_l-V\_dead,0).\par

\textbf{Constraint.} non-draining local lumen volume\par

\textbf{Limit of justification.} An accepted physiological point does not identify this setting or supply a numerical uncertainty interval. No knockout fit or new measurement is performed in this audit.\par

\textbf{Logical domain.} Positive for the retained law; no finite upper physiological bound established.\par

\textbf{Uncertainty status.} genuinely\_unknown\_for\_physiological\_transfer\par

\textbf{Dependence of claims.} Numerical state, current and secretion claims are conditional on this active setting; exact stoichiometric identities are independent of its numerical value.\par

\textbf{Discriminating measurement.} Measure this quantity or its identifiable combination in the same preparation and stimulus protocol, with independent WT data; determine physiological ranges before making global robustness claims.\par

\textbf{Source.} \nolinkurl{src/modern_full_model/water.py:93}. No independently verified primary measurement attached to this numerical setting; model lineage is the source.\par}

\subsubsection*{Record 52: outflow rate s}
{\small \nolinkurl{water.outflow_rate_s}\par

Value: \texttt{0.2}; units: s\textasciicircum{}-1; active: Yes; classification: effective\_assumption.\par

Default: 0.2; differs from default: False.\par

\textbf{Equation role.} Q\_b=L\_b(O\_i-O\_b), Q\_a=L\_a(O\_l-O\_i), Q\_p=L\_p(O\_l-O\_b), Q\_out=k\_out max(V\_l-V\_dead,0).\par

\textbf{Constraint.} local compliance/outflow closure, not secretion calibration\par

\textbf{Limit of justification.} An accepted physiological point does not identify this setting or supply a numerical uncertainty interval. No knockout fit or new measurement is performed in this audit.\par

\textbf{Logical domain.} Positive for the retained law; no finite upper physiological bound established.\par

\textbf{Uncertainty status.} genuinely\_unknown\_for\_physiological\_transfer\par

\textbf{Dependence of claims.} Numerical state, current and secretion claims are conditional on this active setting; exact stoichiometric identities are independent of its numerical value.\par

\textbf{Discriminating measurement.} Measure this quantity or its identifiable combination in the same preparation and stimulus protocol, with independent WT data; determine physiological ranges before making global robustness claims.\par

\textbf{Source.} \nolinkurl{src/modern_full_model/water.py:93}. No independently verified primary measurement attached to this numerical setting; model lineage is the source.\par}

\subsubsection*{Record 53: cell na mM}
{\small \nolinkurl{initial.cell_na_mM}\par

Value: \texttt{15}; units: mM; active: No; classification: constructor\_seed.\par

Default: 15.0; differs from default: False.\par

\textbf{Equation role.} Constructor reference concentrations and volumes define a seed and charge consistency check, not the frozen production initial condition.\par

\textbf{Constraint.} Constructor reference tuple, not the actual production initial state. Task 40 uses the independently frozen WT state; this tuple also defines charge consistency checks.\par

\textbf{Limit of justification.} An accepted physiological point does not identify this setting or supply a numerical uncertainty interval. No knockout fit or new measurement is performed in this audit.\par

\textbf{Logical domain.} Positive for the retained law; no finite upper physiological bound established.\par

\textbf{Uncertainty status.} not\_a\_measured\_uncertainty\_parameter\par

\textbf{Dependence of claims.} No independent parent RHS sensitivity. Used only for the explicitly stated metadata, initial reference, design lineage or inverse family role.\par

\textbf{Discriminating measurement.} No new production parameter experiment is implied for this inactive field; measure the actual corresponding state or active pathway instead.\par

\textbf{Source.} \nolinkurl{src/modern_full_model/model.py:319}. No independently verified primary measurement attached to this numerical setting; model lineage is the source.\par}

\subsubsection*{Record 54: cell k mM}
{\small \nolinkurl{initial.cell_k_mM}\par

Value: \texttt{140}; units: mM; active: No; classification: constructor\_seed.\par

Default: 140.0; differs from default: False.\par

\textbf{Equation role.} Constructor reference concentrations and volumes define a seed and charge consistency check, not the frozen production initial condition.\par

\textbf{Constraint.} Constructor reference tuple, not the actual production initial state. Task 40 uses the independently frozen WT state; this tuple also defines charge consistency checks.\par

\textbf{Limit of justification.} An accepted physiological point does not identify this setting or supply a numerical uncertainty interval. No knockout fit or new measurement is performed in this audit.\par

\textbf{Logical domain.} Positive for the retained law; no finite upper physiological bound established.\par

\textbf{Uncertainty status.} not\_a\_measured\_uncertainty\_parameter\par

\textbf{Dependence of claims.} No independent parent RHS sensitivity. Used only for the explicitly stated metadata, initial reference, design lineage or inverse family role.\par

\textbf{Discriminating measurement.} No new production parameter experiment is implied for this inactive field; measure the actual corresponding state or active pathway instead.\par

\textbf{Source.} \nolinkurl{src/modern_full_model/model.py:319}. No independently verified primary measurement attached to this numerical setting; model lineage is the source.\par}

\subsubsection*{Record 55: cell cl mM}
{\small \nolinkurl{initial.cell_cl_mM}\par

Value: \texttt{40}; units: mM; active: No; classification: constructor\_seed.\par

Default: 40.0; differs from default: False.\par

\textbf{Equation role.} Constructor reference concentrations and volumes define a seed and charge consistency check, not the frozen production initial condition.\par

\textbf{Constraint.} Constructor reference tuple, not the actual production initial state. Task 40 uses the independently frozen WT state; this tuple also defines charge consistency checks.\par

\textbf{Limit of justification.} An accepted physiological point does not identify this setting or supply a numerical uncertainty interval. No knockout fit or new measurement is performed in this audit.\par

\textbf{Logical domain.} Positive for the retained law; no finite upper physiological bound established.\par

\textbf{Uncertainty status.} not\_a\_measured\_uncertainty\_parameter\par

\textbf{Dependence of claims.} No independent parent RHS sensitivity. Used only for the explicitly stated metadata, initial reference, design lineage or inverse family role.\par

\textbf{Discriminating measurement.} No new production parameter experiment is implied for this inactive field; measure the actual corresponding state or active pathway instead.\par

\textbf{Source.} \nolinkurl{src/modern_full_model/model.py:319}. No independently verified primary measurement attached to this numerical setting; model lineage is the source.\par}

\subsubsection*{Record 56: cell tic mM}
{\small \nolinkurl{initial.cell_tic_mM}\par

Value: \texttt{20}; units: mM; active: No; classification: constructor\_seed.\par

Default: 20.0; differs from default: False.\par

\textbf{Equation role.} Constructor reference concentrations and volumes define a seed and charge consistency check, not the frozen production initial condition.\par

\textbf{Constraint.} Constructor reference tuple, not the actual production initial state. Task 40 uses the independently frozen WT state; this tuple also defines charge consistency checks.\par

\textbf{Limit of justification.} An accepted physiological point does not identify this setting or supply a numerical uncertainty interval. No knockout fit or new measurement is performed in this audit.\par

\textbf{Logical domain.} Positive for the retained law; no finite upper physiological bound established.\par

\textbf{Uncertainty status.} not\_a\_measured\_uncertainty\_parameter\par

\textbf{Dependence of claims.} No independent parent RHS sensitivity. Used only for the explicitly stated metadata, initial reference, design lineage or inverse family role.\par

\textbf{Discriminating measurement.} No new production parameter experiment is implied for this inactive field; measure the actual corresponding state or active pathway instead.\par

\textbf{Source.} \nolinkurl{src/modern_full_model/model.py:319}. No independently verified primary measurement attached to this numerical setting; model lineage is the source.\par}

\subsubsection*{Record 57: cell ph}
{\small \nolinkurl{initial.cell_ph}\par

Value: \texttt{7.2}; units: dimensionless (p-scale); active: No; classification: constructor\_seed.\par

Default: 7.2; differs from default: False.\par

\textbf{Equation role.} Constructor reference concentrations and volumes define a seed and charge consistency check, not the frozen production initial condition.\par

\textbf{Constraint.} Constructor reference tuple, not the actual production initial state. Task 40 uses the independently frozen WT state; this tuple also defines charge consistency checks.\par

\textbf{Limit of justification.} An accepted physiological point does not identify this setting or supply a numerical uncertainty interval. No knockout fit or new measurement is performed in this audit.\par

\textbf{Logical domain.} Positive for the retained law; no finite upper physiological bound established.\par

\textbf{Uncertainty status.} not\_a\_measured\_uncertainty\_parameter\par

\textbf{Dependence of claims.} No independent parent RHS sensitivity. Used only for the explicitly stated metadata, initial reference, design lineage or inverse family role.\par

\textbf{Discriminating measurement.} No new production parameter experiment is implied for this inactive field; measure the actual corresponding state or active pathway instead.\par

\textbf{Source.} \nolinkurl{src/modern_full_model/model.py:319}. No independently verified primary measurement attached to this numerical setting; model lineage is the source.\par}

\subsubsection*{Record 58: cell volume pL}
{\small \nolinkurl{initial.cell_volume_pL}\par

Value: \texttt{1}; units: pL; active: No; classification: constructor\_seed.\par

Default: 1.0; differs from default: False.\par

\textbf{Equation role.} Constructor reference concentrations and volumes define a seed and charge consistency check, not the frozen production initial condition.\par

\textbf{Constraint.} Constructor reference tuple, not the actual production initial state. Task 40 uses the independently frozen WT state; this tuple also defines charge consistency checks.\par

\textbf{Limit of justification.} An accepted physiological point does not identify this setting or supply a numerical uncertainty interval. No knockout fit or new measurement is performed in this audit.\par

\textbf{Logical domain.} Positive for the retained law; no finite upper physiological bound established.\par

\textbf{Uncertainty status.} not\_a\_measured\_uncertainty\_parameter\par

\textbf{Dependence of claims.} No independent parent RHS sensitivity. Used only for the explicitly stated metadata, initial reference, design lineage or inverse family role.\par

\textbf{Discriminating measurement.} No new production parameter experiment is implied for this inactive field; measure the actual corresponding state or active pathway instead.\par

\textbf{Source.} \nolinkurl{src/modern_full_model/model.py:319}. No independently verified primary measurement attached to this numerical setting; model lineage is the source.\par}

\subsubsection*{Record 59: lumen na mM}
{\small \nolinkurl{initial.lumen_na_mM}\par

Value: \texttt{145}; units: mM; active: No; classification: constructor\_seed.\par

Default: 145.0; differs from default: False.\par

\textbf{Equation role.} Constructor reference concentrations and volumes define a seed and charge consistency check, not the frozen production initial condition.\par

\textbf{Constraint.} Constructor reference tuple, not the actual production initial state. Task 40 uses the independently frozen WT state; this tuple also defines charge consistency checks.\par

\textbf{Limit of justification.} An accepted physiological point does not identify this setting or supply a numerical uncertainty interval. No knockout fit or new measurement is performed in this audit.\par

\textbf{Logical domain.} Positive for the retained law; no finite upper physiological bound established.\par

\textbf{Uncertainty status.} not\_a\_measured\_uncertainty\_parameter\par

\textbf{Dependence of claims.} No independent parent RHS sensitivity. Used only for the explicitly stated metadata, initial reference, design lineage or inverse family role.\par

\textbf{Discriminating measurement.} No new production parameter experiment is implied for this inactive field; measure the actual corresponding state or active pathway instead.\par

\textbf{Source.} \nolinkurl{src/modern_full_model/model.py:319}. No independently verified primary measurement attached to this numerical setting; model lineage is the source.\par}

\subsubsection*{Record 60: lumen k mM}
{\small \nolinkurl{initial.lumen_k_mM}\par

Value: \texttt{5}; units: mM; active: No; classification: constructor\_seed.\par

Default: 5.0; differs from default: False.\par

\textbf{Equation role.} Constructor reference concentrations and volumes define a seed and charge consistency check, not the frozen production initial condition.\par

\textbf{Constraint.} Constructor reference tuple, not the actual production initial state. Task 40 uses the independently frozen WT state; this tuple also defines charge consistency checks.\par

\textbf{Limit of justification.} An accepted physiological point does not identify this setting or supply a numerical uncertainty interval. No knockout fit or new measurement is performed in this audit.\par

\textbf{Logical domain.} Positive for the retained law; no finite upper physiological bound established.\par

\textbf{Uncertainty status.} not\_a\_measured\_uncertainty\_parameter\par

\textbf{Dependence of claims.} No independent parent RHS sensitivity. Used only for the explicitly stated metadata, initial reference, design lineage or inverse family role.\par

\textbf{Discriminating measurement.} No new production parameter experiment is implied for this inactive field; measure the actual corresponding state or active pathway instead.\par

\textbf{Source.} \nolinkurl{src/modern_full_model/model.py:319}. No independently verified primary measurement attached to this numerical setting; model lineage is the source.\par}

\subsubsection*{Record 61: lumen cl mM}
{\small \nolinkurl{initial.lumen_cl_mM}\par

Value: \texttt{126.1615927965083}; units: mM; active: No; classification: constructor\_seed.\par

Default: 126.16159279650834; differs from default: False.\par

\textbf{Equation role.} Constructor reference concentrations and volumes define a seed and charge consistency check, not the frozen production initial condition.\par

\textbf{Constraint.} Constructor reference tuple, not the actual production initial state. Task 40 uses the independently frozen WT state; this tuple also defines charge consistency checks.\par

\textbf{Limit of justification.} An accepted physiological point does not identify this setting or supply a numerical uncertainty interval. No knockout fit or new measurement is performed in this audit.\par

\textbf{Logical domain.} Positive for the retained law; no finite upper physiological bound established.\par

\textbf{Uncertainty status.} not\_a\_measured\_uncertainty\_parameter\par

\textbf{Dependence of claims.} No independent parent RHS sensitivity. Used only for the explicitly stated metadata, initial reference, design lineage or inverse family role.\par

\textbf{Discriminating measurement.} No new production parameter experiment is implied for this inactive field; measure the actual corresponding state or active pathway instead.\par

\textbf{Source.} \nolinkurl{src/modern_full_model/model.py:319}. No independently verified primary measurement attached to this numerical setting; model lineage is the source.\par}

\subsubsection*{Record 62: lumen tic mM}
{\small \nolinkurl{initial.lumen_tic_mM}\par

Value: \texttt{25}; units: mM; active: No; classification: constructor\_seed.\par

Default: 25.0; differs from default: False.\par

\textbf{Equation role.} Constructor reference concentrations and volumes define a seed and charge consistency check, not the frozen production initial condition.\par

\textbf{Constraint.} Constructor reference tuple, not the actual production initial state. Task 40 uses the independently frozen WT state; this tuple also defines charge consistency checks.\par

\textbf{Limit of justification.} An accepted physiological point does not identify this setting or supply a numerical uncertainty interval. No knockout fit or new measurement is performed in this audit.\par

\textbf{Logical domain.} Positive for the retained law; no finite upper physiological bound established.\par

\textbf{Uncertainty status.} not\_a\_measured\_uncertainty\_parameter\par

\textbf{Dependence of claims.} No independent parent RHS sensitivity. Used only for the explicitly stated metadata, initial reference, design lineage or inverse family role.\par

\textbf{Discriminating measurement.} No new production parameter experiment is implied for this inactive field; measure the actual corresponding state or active pathway instead.\par

\textbf{Source.} \nolinkurl{src/modern_full_model/model.py:319}. No independently verified primary measurement attached to this numerical setting; model lineage is the source.\par}

\subsubsection*{Record 63: lumen ph}
{\small \nolinkurl{initial.lumen_ph}\par

Value: \texttt{7.4}; units: dimensionless (p-scale); active: No; classification: constructor\_seed.\par

Default: 7.4; differs from default: False.\par

\textbf{Equation role.} Constructor reference concentrations and volumes define a seed and charge consistency check, not the frozen production initial condition.\par

\textbf{Constraint.} Constructor reference tuple, not the actual production initial state. Task 40 uses the independently frozen WT state; this tuple also defines charge consistency checks.\par

\textbf{Limit of justification.} An accepted physiological point does not identify this setting or supply a numerical uncertainty interval. No knockout fit or new measurement is performed in this audit.\par

\textbf{Logical domain.} Positive for the retained law; no finite upper physiological bound established.\par

\textbf{Uncertainty status.} not\_a\_measured\_uncertainty\_parameter\par

\textbf{Dependence of claims.} No independent parent RHS sensitivity. Used only for the explicitly stated metadata, initial reference, design lineage or inverse family role.\par

\textbf{Discriminating measurement.} No new production parameter experiment is implied for this inactive field; measure the actual corresponding state or active pathway instead.\par

\textbf{Source.} \nolinkurl{src/modern_full_model/model.py:319}. No independently verified primary measurement attached to this numerical setting; model lineage is the source.\par}

\subsubsection*{Record 64: lumen volume pL}
{\small \nolinkurl{initial.lumen_volume_pL}\par

Value: \texttt{0.1}; units: pL; active: No; classification: constructor\_seed.\par

Default: 0.1; differs from default: False.\par

\textbf{Equation role.} Constructor reference concentrations and volumes define a seed and charge consistency check, not the frozen production initial condition.\par

\textbf{Constraint.} Constructor reference tuple, not the actual production initial state. Task 40 uses the independently frozen WT state; this tuple also defines charge consistency checks.\par

\textbf{Limit of justification.} An accepted physiological point does not identify this setting or supply a numerical uncertainty interval. No knockout fit or new measurement is performed in this audit.\par

\textbf{Logical domain.} Positive for the retained law; no finite upper physiological bound established.\par

\textbf{Uncertainty status.} not\_a\_measured\_uncertainty\_parameter\par

\textbf{Dependence of claims.} No independent parent RHS sensitivity. Used only for the explicitly stated metadata, initial reference, design lineage or inverse family role.\par

\textbf{Discriminating measurement.} No new production parameter experiment is implied for this inactive field; measure the actual corresponding state or active pathway instead.\par

\textbf{Source.} \nolinkurl{src/modern_full_model/model.py:319}. No independently verified primary measurement attached to this numerical setting; model lineage is the source.\par}

\subsubsection*{Record 65: carrier amount fmol}
{\small \nolinkurl{ae4.carrier_amount_fmol}\par

Value: \texttt{0.1501264265070657}; units: fmol; active: Yes; classification: inherited\_effective\_setting.\par

Default: ; differs from default: .\par

\textbf{Equation role.} Shared carrier QSS determines scalar J4; equal routing imposes (-J4/2,-J4/2,J4,-2J4,-2J4) in cell Na,K,Cl,TIC,TA.\par

\textbf{Constraint.} Shared-carrier model setting. Thermodynamic construction fixes forward/reverse ratios, not this absolute abundance or attempt rate.\par

\textbf{Limit of justification.} An accepted physiological point does not identify this setting or supply a numerical uncertainty interval. No knockout fit or new measurement is performed in this audit.\par

\textbf{Logical domain.} Positive for the retained law; no finite upper physiological bound established.\par

\textbf{Uncertainty status.} genuinely\_unknown\_for\_physiological\_transfer\par

\textbf{Dependence of claims.} Numerical state, current and secretion claims are conditional on this active setting; exact stoichiometric identities are independent of its numerical value. Local numerical dependence is quantified in the linked crosswalk; it establishes neither a measured uncertainty range nor global robustness.\par

\textbf{Discriminating measurement.} Measure absolute AE4 abundance and exchange turnover in the same preparation; expression alone does not fix per carrier turnover.\par

\textbf{Source.} \nolinkurl{src/modern_full_model/transporters.py:458}. No independently verified primary measurement attached to this numerical setting; model lineage is the source.\par}

\subsubsection*{Record 66: common cl attempt rate s}
{\small \nolinkurl{ae4.common_cl_attempt_rate_s}\par

Value: \texttt{1}; units: s\textasciicircum{}-1; active: Yes; classification: inherited\_effective\_setting.\par

Default: ; differs from default: .\par

\textbf{Equation role.} Shared carrier QSS determines scalar J4; equal routing imposes (-J4/2,-J4/2,J4,-2J4,-2J4) in cell Na,K,Cl,TIC,TA.\par

\textbf{Constraint.} Shared-carrier model setting. Thermodynamic construction fixes forward/reverse ratios, not this absolute abundance or attempt rate.\par

\textbf{Limit of justification.} An accepted physiological point does not identify this setting or supply a numerical uncertainty interval. No knockout fit or new measurement is performed in this audit.\par

\textbf{Logical domain.} Positive for the retained law; no finite upper physiological bound established.\par

\textbf{Uncertainty status.} genuinely\_unknown\_for\_physiological\_transfer\par

\textbf{Dependence of claims.} Numerical state, current and secretion claims are conditional on this active setting; exact stoichiometric identities are independent of its numerical value.\par

\textbf{Discriminating measurement.} Measure this quantity or its identifiable combination in the same preparation and stimulus protocol, with independent WT data; determine physiological ranges before making global robustness claims.\par

\textbf{Source.} \nolinkurl{src/modern_full_model/transporters.py:458}. No independently verified primary measurement attached to this numerical setting; model lineage is the source.\par}

\subsubsection*{Record 67: na loaded attempt rate s}
{\small \nolinkurl{ae4.na_loaded_attempt_rate_s}\par

Value: \texttt{0.1}; units: s\textasciicircum{}-1; active: Yes; classification: inherited\_effective\_setting.\par

Default: ; differs from default: .\par

\textbf{Equation role.} Shared carrier QSS determines scalar J4; equal routing imposes (-J4/2,-J4/2,J4,-2J4,-2J4) in cell Na,K,Cl,TIC,TA.\par

\textbf{Constraint.} Shared-carrier model setting. Thermodynamic construction fixes forward/reverse ratios, not this absolute abundance or attempt rate.\par

\textbf{Limit of justification.} An accepted physiological point does not identify this setting or supply a numerical uncertainty interval. No knockout fit or new measurement is performed in this audit.\par

\textbf{Logical domain.} Positive for the retained law; no finite upper physiological bound established.\par

\textbf{Uncertainty status.} genuinely\_unknown\_for\_physiological\_transfer\par

\textbf{Dependence of claims.} Numerical state, current and secretion claims are conditional on this active setting; exact stoichiometric identities are independent of its numerical value.\par

\textbf{Discriminating measurement.} Measure this quantity or its identifiable combination in the same preparation and stimulus protocol, with independent WT data; determine physiological ranges before making global robustness claims.\par

\textbf{Source.} \nolinkurl{src/modern_full_model/transporters.py:458}. No independently verified primary measurement attached to this numerical setting; model lineage is the source.\par}

\subsubsection*{Record 68: k loaded attempt rate s}
{\small \nolinkurl{ae4.k_loaded_attempt_rate_s}\par

Value: \texttt{1.9}; units: s\textasciicircum{}-1; active: Yes; classification: inherited\_effective\_setting.\par

Default: ; differs from default: .\par

\textbf{Equation role.} Shared carrier QSS determines scalar J4; equal routing imposes (-J4/2,-J4/2,J4,-2J4,-2J4) in cell Na,K,Cl,TIC,TA.\par

\textbf{Constraint.} Shared-carrier model setting. Thermodynamic construction fixes forward/reverse ratios, not this absolute abundance or attempt rate.\par

\textbf{Limit of justification.} An accepted physiological point does not identify this setting or supply a numerical uncertainty interval. No knockout fit or new measurement is performed in this audit.\par

\textbf{Logical domain.} Positive for the retained law; no finite upper physiological bound established.\par

\textbf{Uncertainty status.} genuinely\_unknown\_for\_physiological\_transfer\par

\textbf{Dependence of claims.} Numerical state, current and secretion claims are conditional on this active setting; exact stoichiometric identities are independent of its numerical value.\par

\textbf{Discriminating measurement.} Measure this quantity or its identifiable combination in the same preparation and stimulus protocol, with independent WT data; determine physiological ranges before making global robustness claims.\par

\textbf{Source.} \nolinkurl{src/modern_full_model/transporters.py:458}. No independently verified primary measurement attached to this numerical setting; model lineage is the source.\par}

\subsubsection*{Record 69: cooperative gate}
{\small \nolinkurl{ae4.cooperative_gate}\par

Value: None; units: configuration; active: No; classification: inherited\_effective\_setting.\par

Default: ; differs from default: .\par

\textbf{Equation role.} Shared carrier QSS determines scalar J4; equal routing imposes (-J4/2,-J4/2,J4,-2J4,-2J4) in cell Na,K,Cl,TIC,TA.\par

\textbf{Constraint.} Shared-carrier model setting. Thermodynamic construction fixes forward/reverse ratios, not this absolute abundance or attempt rate.\par

\textbf{Limit of justification.} An accepted physiological point does not identify this setting or supply a numerical uncertainty interval. No knockout fit or new measurement is performed in this audit.\par

\textbf{Logical domain.} Selected discrete setting or stated numerical reference.\par

\textbf{Uncertainty status.} genuinely\_unknown\_for\_physiological\_transfer\par

\textbf{Dependence of claims.} No independent parent RHS sensitivity. Used only for the explicitly stated metadata, initial reference, design lineage or inverse family role.\par

\textbf{Discriminating measurement.} Measure this quantity or its identifiable combination in the same preparation and stimulus protocol, with independent WT data; determine physiological ranges before making global robustness claims.\par

\textbf{Source.} \nolinkurl{src/modern_full_model/transporters.py:458}. No independently verified primary measurement attached to this numerical setting; model lineage is the source.\par}

\subsubsection*{Record 70: capacity fmol s}
{\small \nolinkurl{nbc.capacity_fmol_s}\par

Value: \texttt{0.115709131974644}; units: fmol/s; active: Yes; classification: WT\_architecture\_derived.\par

Default: ; differs from default: .\par

\textbf{Equation role.} JB=G\_B u(Ca) tanh(affinity/log\_width), with electrogenic 1 Na and 2 bicarbonate inward per cycle.\par

\textbf{Constraint.} Derived from an assumed 70/30 positive chloride loading architecture at the Task 31 reference state, a 2.3 NHE1 multiplier and coupled electrical closure. Derived conditional on those assumptions, not measured or identified by knockout data.\par

\textbf{Limit of justification.} The 70/30 partition is assumed, not measured. Solving the coupled current balance does not establish the physiological truth or unique identification of that partition.\par

\textbf{Logical domain.} Positive for the retained law; no finite upper physiological bound established.\par

\textbf{Uncertainty status.} conditional\_derivation\_no\_experimental\_interval\par

\textbf{Dependence of claims.} Numerical state, current and secretion claims are conditional on this active setting; exact stoichiometric identities are independent of its numerical value. Local numerical dependence is quantified in the linked crosswalk; it establishes neither a measured uncertainty range nor global robustness.\par

\textbf{Discriminating measurement.} Measure Na dependent bicarbonate loading or alkalinity recovery during stimulation with calibrated pH, TIC and volume, separating NHE1 and AE2 contributions; establish molecular identity independently.\par

\textbf{Source.} \nolinkurl{src/modern_full_model/nbc_minimal.py:190}. No independently verified primary measurement attached to this numerical setting; model lineage is the source.\par}

\subsubsection*{Record 71: log width}
{\small \nolinkurl{nbc.log_width}\par

Value: \texttt{2}; units: dimensionless; active: Yes; classification: effective\_assumption.\par

Default: ; differs from default: .\par

\textbf{Equation role.} JB=G\_B u(Ca) tanh(affinity/log\_width), with electrogenic 1 Na and 2 bicarbonate inward per cycle.\par

\textbf{Constraint.} Inherited NBC recruitment parameter; not an isoform-specific measurement.\par

\textbf{Limit of justification.} An accepted physiological point does not identify this setting or supply a numerical uncertainty interval. No knockout fit or new measurement is performed in this audit.\par

\textbf{Logical domain.} Positive for the retained law; no finite upper physiological bound established.\par

\textbf{Uncertainty status.} genuinely\_unknown\_for\_physiological\_transfer\par

\textbf{Dependence of claims.} Numerical state, current and secretion claims are conditional on this active setting; exact stoichiometric identities are independent of its numerical value. Local numerical dependence is quantified in the linked crosswalk; it establishes neither a measured uncertainty range nor global robustness.\par

\textbf{Discriminating measurement.} Measure this quantity or its identifiable combination in the same preparation and stimulus protocol, with independent WT data; determine physiological ranges before making global robustness claims.\par

\textbf{Source.} \nolinkurl{src/modern_full_model/nbc_minimal.py:190}. No independently verified primary measurement attached to this numerical setting; model lineage is the source.\par}

\subsubsection*{Record 72: resting calcium uM}
{\small \nolinkurl{nbc.resting_calcium_uM}\par

Value: \texttt{0.058}; units: uM; active: Yes; classification: effective\_assumption.\par

Default: ; differs from default: .\par

\textbf{Equation role.} JB=G\_B u(Ca) tanh(affinity/log\_width), with electrogenic 1 Na and 2 bicarbonate inward per cycle.\par

\textbf{Constraint.} Inherited NBC recruitment parameter; not an isoform-specific measurement.\par

\textbf{Limit of justification.} An accepted physiological point does not identify this setting or supply a numerical uncertainty interval. No knockout fit or new measurement is performed in this audit.\par

\textbf{Logical domain.} Positive for the retained law; no finite upper physiological bound established.\par

\textbf{Uncertainty status.} genuinely\_unknown\_for\_physiological\_transfer\par

\textbf{Dependence of claims.} Numerical state, current and secretion claims are conditional on this active setting; exact stoichiometric identities are independent of its numerical value.\par

\textbf{Discriminating measurement.} Measure this quantity or its identifiable combination in the same preparation and stimulus protocol, with independent WT data; determine physiological ranges before making global robustness claims.\par

\textbf{Source.} \nolinkurl{src/modern_full_model/nbc_minimal.py:190}. No independently verified primary measurement attached to this numerical setting; model lineage is the source.\par}

\subsubsection*{Record 73: fully recruited calcium uM}
{\small \nolinkurl{nbc.fully_recruited_calcium_uM}\par

Value: \texttt{0.25}; units: uM; active: Yes; classification: effective\_assumption.\par

Default: ; differs from default: .\par

\textbf{Equation role.} JB=G\_B u(Ca) tanh(affinity/log\_width), with electrogenic 1 Na and 2 bicarbonate inward per cycle.\par

\textbf{Constraint.} Inherited NBC recruitment parameter; not an isoform-specific measurement.\par

\textbf{Limit of justification.} An accepted physiological point does not identify this setting or supply a numerical uncertainty interval. No knockout fit or new measurement is performed in this audit.\par

\textbf{Logical domain.} Positive for the retained law; no finite upper physiological bound established.\par

\textbf{Uncertainty status.} genuinely\_unknown\_for\_physiological\_transfer\par

\textbf{Dependence of claims.} Numerical state, current and secretion claims are conditional on this active setting; exact stoichiometric identities are independent of its numerical value.\par

\textbf{Discriminating measurement.} Measure this quantity or its identifiable combination in the same preparation and stimulus protocol, with independent WT data; determine physiological ranges before making global robustness claims.\par

\textbf{Source.} \nolinkurl{src/modern_full_model/nbc_minimal.py:190}. No independently verified primary measurement attached to this numerical setting; model lineage is the source.\par}

\subsubsection*{Record 74: nhe1 stimulated multiplier}
{\small \nolinkurl{nbc.nhe1_stimulated_multiplier}\par

Value: \texttt{2.3}; units: dimensionless; active: Yes; classification: effective\_assumption.\par

Default: ; differs from default: .\par

\textbf{Equation role.} JH\_stim=[1+(multiplier-1)*u(Ca)]*JH\_Cha; shares the NBC recruitment endpoints.\par

\textbf{Constraint.} Inherited source-fixed 2.3 multiplier. Fixing it before a test prevents retuning but does not make its transfer to this tissue/protocol uncertainty-free.\par

\textbf{Limit of justification.} An accepted physiological point does not identify this setting or supply a numerical uncertainty interval. No knockout fit or new measurement is performed in this audit.\par

\textbf{Logical domain.} At least 1 by the selected model constructor; this is a structural choice, not an experimental lower bound.\par

\textbf{Uncertainty status.} genuinely\_unknown\_for\_physiological\_transfer\par

\textbf{Dependence of claims.} Numerical state, current and secretion claims are conditional on this active setting; exact stoichiometric identities are independent of its numerical value. Local numerical dependence is quantified in the linked crosswalk; it establishes neither a measured uncertainty range nor global robustness.\par

\textbf{Discriminating measurement.} Measure the ratio of stimulated to resting NHE1 flux in the same salivary preparation and calcium protocol.\par

\textbf{Source.} \nolinkurl{src/modern_full_model/nbc_minimal.py:149}. No independently verified primary measurement attached to this numerical setting; model lineage is the source.\par}

\subsubsection*{Record 75: alpha eff fmol s}
{\small \nolinkurl{nkcc.alpha_eff_fmol_s}\par

Value: \texttt{0.01733474689409605}; units: fmol/s; active: Yes; classification: WT\_constraint\_derived.\par

Default: ; differs from default: .\par

\textbf{Equation role.} JN=alpha*m\_N*(A1-A2[Na]i[K]i[Cl]i\textasciicircum{}2)/(A3+A4[Na]i[K]i[Cl]i\textasciicircum{}2); inward 1 Na, 1 K, 2 Cl per cycle.\par

\textbf{Constraint.} J\_N,rest divided by the published shape at the accepted resting concentrations. The resting flux is itself an inherited model value, not a direct new measurement.\par

\textbf{Limit of justification.} An accepted physiological point does not identify this setting or supply a numerical uncertainty interval. No knockout fit or new measurement is performed in this audit.\par

\textbf{Logical domain.} Positive for the retained law; no finite upper physiological bound established.\par

\textbf{Uncertainty status.} conditional\_derivation\_no\_experimental\_interval\par

\textbf{Dependence of claims.} Numerical state, current and secretion claims are conditional on this active setting; exact stoichiometric identities are independent of its numerical value. Local numerical dependence is quantified in the linked crosswalk; it establishes neither a measured uncertainty range nor global robustness.\par

\textbf{Discriminating measurement.} Measure acute intact cell NKCC chloride uptake during AE4 loss with simultaneous intracellular Na, K and Cl; compare separately with the isolated 2015 uptake assay.\par

\textbf{Source.} \nolinkurl{src/modern_full_model/nkcc1_palk2010.py:31}. Palk et al. 2010/Benjamin and Johnson 1997 law, reproduced as Vera Siguenza et al. 2018 Eq. 17; coefficients in src/modern\_full\_model/nkcc1\_palk2010.py.\par}

\subsubsection*{Record 76: A1}
{\small \nolinkurl{nkcc.A1}\par

Value: \texttt{157.5}; units: dimensionless; active: Yes; classification: published\_kinetic\_coefficient.\par

Default: ; differs from default: .\par

\textbf{Equation role.} JN=alpha*m\_N*(A1-A2[Na]i[K]i[Cl]i\textasciicircum{}2)/(A3+A4[Na]i[K]i[Cl]i\textasciicircum{}2); inward 1 Na, 1 K, 2 Cl per cycle.\par

\textbf{Constraint.} Palk/Benjamin effective two-state law in the fixed bath, with fourth-order concentration factors converted to mM. Do not vary bath composition without revisiting these fixed-bath coefficients.\par

\textbf{Limit of justification.} Publication of the kinetic law does not measure salivary carrier abundance or establish an uncertainty interval for transfer to this preparation.\par

\textbf{Logical domain.} Positive for the retained law; no finite upper physiological bound established.\par

\textbf{Uncertainty status.} genuinely\_unknown\_for\_physiological\_transfer\par

\textbf{Dependence of claims.} Numerical state, current and secretion claims are conditional on this active setting; exact stoichiometric identities are independent of its numerical value.\par

\textbf{Discriminating measurement.} Measure this quantity or its identifiable combination in the same preparation and stimulus protocol, with independent WT data; determine physiological ranges before making global robustness claims.\par

\textbf{Source.} \nolinkurl{src/modern_full_model/nkcc1_palk2010.py:31}. Palk et al. 2010/Benjamin and Johnson 1997 law, reproduced as Vera Siguenza et al. 2018 Eq. 17; coefficients in src/modern\_full\_model/nkcc1\_palk2010.py.\par}

\subsubsection*{Record 77: A2}
{\small \nolinkurl{nkcc.A2}\par

Value: \texttt{2.0096e-05}; units: mM\textasciicircum{}-4; active: Yes; classification: published\_kinetic\_coefficient.\par

Default: ; differs from default: .\par

\textbf{Equation role.} JN=alpha*m\_N*(A1-A2[Na]i[K]i[Cl]i\textasciicircum{}2)/(A3+A4[Na]i[K]i[Cl]i\textasciicircum{}2); inward 1 Na, 1 K, 2 Cl per cycle.\par

\textbf{Constraint.} Palk/Benjamin effective two-state law in the fixed bath, with fourth-order concentration factors converted to mM. Do not vary bath composition without revisiting these fixed-bath coefficients.\par

\textbf{Limit of justification.} Publication of the kinetic law does not measure salivary carrier abundance or establish an uncertainty interval for transfer to this preparation.\par

\textbf{Logical domain.} Positive for the retained law; no finite upper physiological bound established.\par

\textbf{Uncertainty status.} genuinely\_unknown\_for\_physiological\_transfer\par

\textbf{Dependence of claims.} Numerical state, current and secretion claims are conditional on this active setting; exact stoichiometric identities are independent of its numerical value.\par

\textbf{Discriminating measurement.} Measure this quantity or its identifiable combination in the same preparation and stimulus protocol, with independent WT data; determine physiological ranges before making global robustness claims.\par

\textbf{Source.} \nolinkurl{src/modern_full_model/nkcc1_palk2010.py:31}. Palk et al. 2010/Benjamin and Johnson 1997 law, reproduced as Vera Siguenza et al. 2018 Eq. 17; coefficients in src/modern\_full\_model/nkcc1\_palk2010.py.\par}

\subsubsection*{Record 78: A3}
{\small \nolinkurl{nkcc.A3}\par

Value: \texttt{1.0306}; units: dimensionless; active: Yes; classification: published\_kinetic\_coefficient.\par

Default: ; differs from default: .\par

\textbf{Equation role.} JN=alpha*m\_N*(A1-A2[Na]i[K]i[Cl]i\textasciicircum{}2)/(A3+A4[Na]i[K]i[Cl]i\textasciicircum{}2); inward 1 Na, 1 K, 2 Cl per cycle.\par

\textbf{Constraint.} Palk/Benjamin effective two-state law in the fixed bath, with fourth-order concentration factors converted to mM. Do not vary bath composition without revisiting these fixed-bath coefficients.\par

\textbf{Limit of justification.} Publication of the kinetic law does not measure salivary carrier abundance or establish an uncertainty interval for transfer to this preparation.\par

\textbf{Logical domain.} Positive for the retained law; no finite upper physiological bound established.\par

\textbf{Uncertainty status.} genuinely\_unknown\_for\_physiological\_transfer\par

\textbf{Dependence of claims.} Numerical state, current and secretion claims are conditional on this active setting; exact stoichiometric identities are independent of its numerical value.\par

\textbf{Discriminating measurement.} Measure this quantity or its identifiable combination in the same preparation and stimulus protocol, with independent WT data; determine physiological ranges before making global robustness claims.\par

\textbf{Source.} \nolinkurl{src/modern_full_model/nkcc1_palk2010.py:31}. Palk et al. 2010/Benjamin and Johnson 1997 law, reproduced as Vera Siguenza et al. 2018 Eq. 17; coefficients in src/modern\_full\_model/nkcc1\_palk2010.py.\par}

\subsubsection*{Record 79: A4}
{\small \nolinkurl{nkcc.A4}\par

Value: \texttt{1.3852e-06}; units: mM\textasciicircum{}-4; active: Yes; classification: published\_kinetic\_coefficient.\par

Default: ; differs from default: .\par

\textbf{Equation role.} JN=alpha*m\_N*(A1-A2[Na]i[K]i[Cl]i\textasciicircum{}2)/(A3+A4[Na]i[K]i[Cl]i\textasciicircum{}2); inward 1 Na, 1 K, 2 Cl per cycle.\par

\textbf{Constraint.} Palk/Benjamin effective two-state law in the fixed bath, with fourth-order concentration factors converted to mM. Do not vary bath composition without revisiting these fixed-bath coefficients.\par

\textbf{Limit of justification.} Publication of the kinetic law does not measure salivary carrier abundance or establish an uncertainty interval for transfer to this preparation.\par

\textbf{Logical domain.} Positive for the retained law; no finite upper physiological bound established.\par

\textbf{Uncertainty status.} genuinely\_unknown\_for\_physiological\_transfer\par

\textbf{Dependence of claims.} Numerical state, current and secretion claims are conditional on this active setting; exact stoichiometric identities are independent of its numerical value.\par

\textbf{Discriminating measurement.} Measure this quantity or its identifiable combination in the same preparation and stimulus protocol, with independent WT data; determine physiological ranges before making global robustness claims.\par

\textbf{Source.} \nolinkurl{src/modern_full_model/nkcc1_palk2010.py:31}. Palk et al. 2010/Benjamin and Johnson 1997 law, reproduced as Vera Siguenza et al. 2018 Eq. 17; coefficients in src/modern\_full\_model/nkcc1\_palk2010.py.\par}

\subsubsection*{Record 80: k1 plus per ms}
{\small \nolinkurl{cha.k1_plus_per_ms}\par

Value: \texttt{10.5}; units: ms\textasciicircum{}-1; active: Yes; classification: published\_kinetic\_coefficient.\par

Default: ; differs from default: .\par

\textbf{Equation role.} JH=N\_H*m\_H*Mod2*(a*b-c*d)/(a+b+c+d), with rate conversion from ms to s and printed fourth reverse rate retained.\par

\textbf{Constraint.} Printed Cha 2009 Table S1 coefficient, transferred as a kinetic law. The printed fourth reverse rate is retained despite its documented rounding discrepancy from exact microscopic reversibility.\par

\textbf{Limit of justification.} Publication of the kinetic law does not measure salivary carrier abundance or establish an uncertainty interval for transfer to this preparation.\par

\textbf{Logical domain.} Positive for the retained law; no finite upper physiological bound established.\par

\textbf{Uncertainty status.} genuinely\_unknown\_for\_physiological\_transfer\par

\textbf{Dependence of claims.} Numerical state, current and secretion claims are conditional on this active setting; exact stoichiometric identities are independent of its numerical value.\par

\textbf{Discriminating measurement.} Measure this quantity or its identifiable combination in the same preparation and stimulus protocol, with independent WT data; determine physiological ranges before making global robustness claims.\par

\textbf{Source.} \nolinkurl{src/modern_full_model/nhe1_cha2009.py:113}. Cha et al. 2009, Eqs. 3 and 5, Table S1; DOI 10.1016/j.bpj.2009.08.053; repository results/31\_nhe1\_mechanistic\_repair/source\_table\_s1.json.\par}

\subsubsection*{Record 81: k1 minus per ms}
{\small \nolinkurl{cha.k1_minus_per_ms}\par

Value: \texttt{0.201}; units: ms\textasciicircum{}-1; active: Yes; classification: published\_kinetic\_coefficient.\par

Default: ; differs from default: .\par

\textbf{Equation role.} JH=N\_H*m\_H*Mod2*(a*b-c*d)/(a+b+c+d), with rate conversion from ms to s and printed fourth reverse rate retained.\par

\textbf{Constraint.} Printed Cha 2009 Table S1 coefficient, transferred as a kinetic law. The printed fourth reverse rate is retained despite its documented rounding discrepancy from exact microscopic reversibility.\par

\textbf{Limit of justification.} Publication of the kinetic law does not measure salivary carrier abundance or establish an uncertainty interval for transfer to this preparation.\par

\textbf{Logical domain.} Positive for the retained law; no finite upper physiological bound established.\par

\textbf{Uncertainty status.} genuinely\_unknown\_for\_physiological\_transfer\par

\textbf{Dependence of claims.} Numerical state, current and secretion claims are conditional on this active setting; exact stoichiometric identities are independent of its numerical value.\par

\textbf{Discriminating measurement.} Measure this quantity or its identifiable combination in the same preparation and stimulus protocol, with independent WT data; determine physiological ranges before making global robustness claims.\par

\textbf{Source.} \nolinkurl{src/modern_full_model/nhe1_cha2009.py:113}. Cha et al. 2009, Eqs. 3 and 5, Table S1; DOI 10.1016/j.bpj.2009.08.053; repository results/31\_nhe1\_mechanistic\_repair/source\_table\_s1.json.\par}

\subsubsection*{Record 82: k2 plus per ms}
{\small \nolinkurl{cha.k2_plus_per_ms}\par

Value: \texttt{15.8}; units: ms\textasciicircum{}-1; active: Yes; classification: published\_kinetic\_coefficient.\par

Default: ; differs from default: .\par

\textbf{Equation role.} JH=N\_H*m\_H*Mod2*(a*b-c*d)/(a+b+c+d), with rate conversion from ms to s and printed fourth reverse rate retained.\par

\textbf{Constraint.} Printed Cha 2009 Table S1 coefficient, transferred as a kinetic law. The printed fourth reverse rate is retained despite its documented rounding discrepancy from exact microscopic reversibility.\par

\textbf{Limit of justification.} Publication of the kinetic law does not measure salivary carrier abundance or establish an uncertainty interval for transfer to this preparation.\par

\textbf{Logical domain.} Positive for the retained law; no finite upper physiological bound established.\par

\textbf{Uncertainty status.} genuinely\_unknown\_for\_physiological\_transfer\par

\textbf{Dependence of claims.} Numerical state, current and secretion claims are conditional on this active setting; exact stoichiometric identities are independent of its numerical value.\par

\textbf{Discriminating measurement.} Measure this quantity or its identifiable combination in the same preparation and stimulus protocol, with independent WT data; determine physiological ranges before making global robustness claims.\par

\textbf{Source.} \nolinkurl{src/modern_full_model/nhe1_cha2009.py:113}. Cha et al. 2009, Eqs. 3 and 5, Table S1; DOI 10.1016/j.bpj.2009.08.053; repository results/31\_nhe1\_mechanistic\_repair/source\_table\_s1.json.\par}

\subsubsection*{Record 83: k na o mM}
{\small \nolinkurl{cha.k_na_o_mM}\par

Value: \texttt{195}; units: mM; active: Yes; classification: published\_kinetic\_coefficient.\par

Default: ; differs from default: .\par

\textbf{Equation role.} JH=N\_H*m\_H*Mod2*(a*b-c*d)/(a+b+c+d), with rate conversion from ms to s and printed fourth reverse rate retained.\par

\textbf{Constraint.} Printed Cha 2009 Table S1 coefficient, transferred as a kinetic law. The printed fourth reverse rate is retained despite its documented rounding discrepancy from exact microscopic reversibility.\par

\textbf{Limit of justification.} Publication of the kinetic law does not measure salivary carrier abundance or establish an uncertainty interval for transfer to this preparation.\par

\textbf{Logical domain.} Positive for the retained law; no finite upper physiological bound established.\par

\textbf{Uncertainty status.} genuinely\_unknown\_for\_physiological\_transfer\par

\textbf{Dependence of claims.} Numerical state, current and secretion claims are conditional on this active setting; exact stoichiometric identities are independent of its numerical value.\par

\textbf{Discriminating measurement.} Measure this quantity or its identifiable combination in the same preparation and stimulus protocol, with independent WT data; determine physiological ranges before making global robustness claims.\par

\textbf{Source.} \nolinkurl{src/modern_full_model/nhe1_cha2009.py:113}. Cha et al. 2009, Eqs. 3 and 5, Table S1; DOI 10.1016/j.bpj.2009.08.053; repository results/31\_nhe1\_mechanistic\_repair/source\_table\_s1.json.\par}

\subsubsection*{Record 84: k na i mM}
{\small \nolinkurl{cha.k_na_i_mM}\par

Value: \texttt{16.2}; units: mM; active: Yes; classification: published\_kinetic\_coefficient.\par

Default: ; differs from default: .\par

\textbf{Equation role.} JH=N\_H*m\_H*Mod2*(a*b-c*d)/(a+b+c+d), with rate conversion from ms to s and printed fourth reverse rate retained.\par

\textbf{Constraint.} Printed Cha 2009 Table S1 coefficient, transferred as a kinetic law. The printed fourth reverse rate is retained despite its documented rounding discrepancy from exact microscopic reversibility.\par

\textbf{Limit of justification.} Publication of the kinetic law does not measure salivary carrier abundance or establish an uncertainty interval for transfer to this preparation.\par

\textbf{Logical domain.} Positive for the retained law; no finite upper physiological bound established.\par

\textbf{Uncertainty status.} genuinely\_unknown\_for\_physiological\_transfer\par

\textbf{Dependence of claims.} Numerical state, current and secretion claims are conditional on this active setting; exact stoichiometric identities are independent of its numerical value.\par

\textbf{Discriminating measurement.} Measure this quantity or its identifiable combination in the same preparation and stimulus protocol, with independent WT data; determine physiological ranges before making global robustness claims.\par

\textbf{Source.} \nolinkurl{src/modern_full_model/nhe1_cha2009.py:113}. Cha et al. 2009, Eqs. 3 and 5, Table S1; DOI 10.1016/j.bpj.2009.08.053; repository results/31\_nhe1\_mechanistic\_repair/source\_table\_s1.json.\par}

\subsubsection*{Record 85: k h o mM}
{\small \nolinkurl{cha.k_h_o_mM}\par

Value: \texttt{0.00162}; units: mM; active: Yes; classification: published\_kinetic\_coefficient.\par

Default: ; differs from default: .\par

\textbf{Equation role.} JH=N\_H*m\_H*Mod2*(a*b-c*d)/(a+b+c+d), with rate conversion from ms to s and printed fourth reverse rate retained.\par

\textbf{Constraint.} Printed Cha 2009 Table S1 coefficient, transferred as a kinetic law. The printed fourth reverse rate is retained despite its documented rounding discrepancy from exact microscopic reversibility.\par

\textbf{Limit of justification.} Publication of the kinetic law does not measure salivary carrier abundance or establish an uncertainty interval for transfer to this preparation.\par

\textbf{Logical domain.} Positive for the retained law; no finite upper physiological bound established.\par

\textbf{Uncertainty status.} genuinely\_unknown\_for\_physiological\_transfer\par

\textbf{Dependence of claims.} Numerical state, current and secretion claims are conditional on this active setting; exact stoichiometric identities are independent of its numerical value.\par

\textbf{Discriminating measurement.} Measure this quantity or its identifiable combination in the same preparation and stimulus protocol, with independent WT data; determine physiological ranges before making global robustness claims.\par

\textbf{Source.} \nolinkurl{src/modern_full_model/nhe1_cha2009.py:113}. Cha et al. 2009, Eqs. 3 and 5, Table S1; DOI 10.1016/j.bpj.2009.08.053; repository results/31\_nhe1\_mechanistic\_repair/source\_table\_s1.json.\par}

\subsubsection*{Record 86: k h i mM}
{\small \nolinkurl{cha.k_h_i_mM}\par

Value: \texttt{0.000605}; units: mM; active: Yes; classification: published\_kinetic\_coefficient.\par

Default: ; differs from default: .\par

\textbf{Equation role.} JH=N\_H*m\_H*Mod2*(a*b-c*d)/(a+b+c+d), with rate conversion from ms to s and printed fourth reverse rate retained.\par

\textbf{Constraint.} Printed Cha 2009 Table S1 coefficient, transferred as a kinetic law. The printed fourth reverse rate is retained despite its documented rounding discrepancy from exact microscopic reversibility.\par

\textbf{Limit of justification.} Publication of the kinetic law does not measure salivary carrier abundance or establish an uncertainty interval for transfer to this preparation.\par

\textbf{Logical domain.} Positive for the retained law; no finite upper physiological bound established.\par

\textbf{Uncertainty status.} genuinely\_unknown\_for\_physiological\_transfer\par

\textbf{Dependence of claims.} Numerical state, current and secretion claims are conditional on this active setting; exact stoichiometric identities are independent of its numerical value.\par

\textbf{Discriminating measurement.} Measure this quantity or its identifiable combination in the same preparation and stimulus protocol, with independent WT data; determine physiological ranges before making global robustness claims.\par

\textbf{Source.} \nolinkurl{src/modern_full_model/nhe1_cha2009.py:113}. Cha et al. 2009, Eqs. 3 and 5, Table S1; DOI 10.1016/j.bpj.2009.08.053; repository results/31\_nhe1\_mechanistic\_repair/source\_table\_s1.json.\par}

\subsubsection*{Record 87: modifier k i mM}
{\small \nolinkurl{cha.modifier_k_i_mM}\par

Value: \texttt{3.07e-05}; units: mM; active: Yes; classification: published\_kinetic\_coefficient.\par

Default: ; differs from default: .\par

\textbf{Equation role.} JH=N\_H*m\_H*Mod2*(a*b-c*d)/(a+b+c+d), with rate conversion from ms to s and printed fourth reverse rate retained.\par

\textbf{Constraint.} Printed Cha 2009 Table S1 coefficient, transferred as a kinetic law. The printed fourth reverse rate is retained despite its documented rounding discrepancy from exact microscopic reversibility.\par

\textbf{Limit of justification.} Publication of the kinetic law does not measure salivary carrier abundance or establish an uncertainty interval for transfer to this preparation.\par

\textbf{Logical domain.} Positive for the retained law; no finite upper physiological bound established.\par

\textbf{Uncertainty status.} genuinely\_unknown\_for\_physiological\_transfer\par

\textbf{Dependence of claims.} Numerical state, current and secretion claims are conditional on this active setting; exact stoichiometric identities are independent of its numerical value.\par

\textbf{Discriminating measurement.} Measure this quantity or its identifiable combination in the same preparation and stimulus protocol, with independent WT data; determine physiological ranges before making global robustness claims.\par

\textbf{Source.} \nolinkurl{src/modern_full_model/nhe1_cha2009.py:113}. Cha et al. 2009, Eqs. 3 and 5, Table S1; DOI 10.1016/j.bpj.2009.08.053; repository results/31\_nhe1\_mechanistic\_repair/source\_table\_s1.json.\par}

\subsubsection*{Record 88: modifier k o mM}
{\small \nolinkurl{cha.modifier_k_o_mM}\par

Value: \texttt{4.8e-07}; units: mM; active: Yes; classification: published\_kinetic\_coefficient.\par

Default: ; differs from default: .\par

\textbf{Equation role.} JH=N\_H*m\_H*Mod2*(a*b-c*d)/(a+b+c+d), with rate conversion from ms to s and printed fourth reverse rate retained.\par

\textbf{Constraint.} Printed Cha 2009 Table S1 coefficient, transferred as a kinetic law. The printed fourth reverse rate is retained despite its documented rounding discrepancy from exact microscopic reversibility.\par

\textbf{Limit of justification.} Publication of the kinetic law does not measure salivary carrier abundance or establish an uncertainty interval for transfer to this preparation.\par

\textbf{Logical domain.} Positive for the retained law; no finite upper physiological bound established.\par

\textbf{Uncertainty status.} genuinely\_unknown\_for\_physiological\_transfer\par

\textbf{Dependence of claims.} Numerical state, current and secretion claims are conditional on this active setting; exact stoichiometric identities are independent of its numerical value.\par

\textbf{Discriminating measurement.} Measure this quantity or its identifiable combination in the same preparation and stimulus protocol, with independent WT data; determine physiological ranges before making global robustness claims.\par

\textbf{Source.} \nolinkurl{src/modern_full_model/nhe1_cha2009.py:113}. Cha et al. 2009, Eqs. 3 and 5, Table S1; DOI 10.1016/j.bpj.2009.08.053; repository results/31\_nhe1\_mechanistic\_repair/source\_table\_s1.json.\par}

\subsubsection*{Record 89: k2 minus table per ms}
{\small \nolinkurl{cha.k2_minus_table_per_ms}\par

Value: \texttt{183}; units: ms\textasciicircum{}-1; active: Yes; classification: published\_kinetic\_coefficient.\par

Default: ; differs from default: .\par

\textbf{Equation role.} JH=N\_H*m\_H*Mod2*(a*b-c*d)/(a+b+c+d), with rate conversion from ms to s and printed fourth reverse rate retained.\par

\textbf{Constraint.} Printed Cha 2009 Table S1 coefficient, transferred as a kinetic law. The printed fourth reverse rate is retained despite its documented rounding discrepancy from exact microscopic reversibility.\par

\textbf{Limit of justification.} Publication of the kinetic law does not measure salivary carrier abundance or establish an uncertainty interval for transfer to this preparation.\par

\textbf{Logical domain.} Positive for the retained law; no finite upper physiological bound established.\par

\textbf{Uncertainty status.} genuinely\_unknown\_for\_physiological\_transfer\par

\textbf{Dependence of claims.} Numerical state, current and secretion claims are conditional on this active setting; exact stoichiometric identities are independent of its numerical value.\par

\textbf{Discriminating measurement.} Measure this quantity or its identifiable combination in the same preparation and stimulus protocol, with independent WT data; determine physiological ranges before making global robustness claims.\par

\textbf{Source.} \nolinkurl{src/modern_full_model/nhe1_cha2009.py:113}. Cha et al. 2009, Eqs. 3 and 5, Table S1; DOI 10.1016/j.bpj.2009.08.053; repository results/31\_nhe1\_mechanistic\_repair/source\_table\_s1.json.\par}

\subsubsection*{Record 90: tau ae4 s}
{\small \nolinkurl{regulation.tau_ae4_s}\par

Value: \texttt{30}; units: s; active: Yes; classification: effective\_assumption.\par

Default: ; differs from default: .\par

\textbf{Equation role.} da/dt=(beta-a)/tau; AE4 multiplier=basal+increment*coupling*construct\_scale*a.\par

\textbf{Constraint.} A 30 s effective response time, not a resolved measurement of the intermediate signalling kinetics.\par

\textbf{Limit of justification.} An accepted physiological point does not identify this setting or supply a numerical uncertainty interval. No knockout fit or new measurement is performed in this audit.\par

\textbf{Logical domain.} Positive for the retained law; no finite upper physiological bound established.\par

\textbf{Uncertainty status.} genuinely\_unknown\_for\_physiological\_transfer\par

\textbf{Dependence of claims.} Regulatory transient claims depend on the selected family/timescale and gain; constant full stimulus equilibrium is independent of tau within R1. Local numerical dependence is quantified in the linked crosswalk; it establishes neither a measured uncertainty range nor global robustness.\par

\textbf{Discriminating measurement.} Acquire acute beta input, cAMP/PKA activity and AE4 transport time series to distinguish regulation from slow ionic and volume modes.\par

\textbf{Source.} \nolinkurl{src/modern_full_model/camp_pka.py:90}. Pena Munzenmayer et al. 2021 gives pathway and endpoint activation evidence; acute kinetic times and exact common capacity gain remain effective settings.\par}

\subsubsection*{Record 91: basal capacity multiplier}
{\small \nolinkurl{regulation.basal_capacity_multiplier}\par

Value: \texttt{1}; units: dimensionless; active: Yes; classification: effective\_assumption.\par

Default: ; differs from default: .\par

\textbf{Equation role.} da/dt=(beta-a)/tau; AE4 multiplier=basal+increment*coupling*construct\_scale*a.\par

\textbf{Constraint.} Effective normalisation or gain; pathway activation evidence does not identify each coefficient independently.\par

\textbf{Limit of justification.} An accepted physiological point does not identify this setting or supply a numerical uncertainty interval. No knockout fit or new measurement is performed in this audit.\par

\textbf{Logical domain.} Positive for the retained law; no finite upper physiological bound established.\par

\textbf{Uncertainty status.} genuinely\_unknown\_for\_physiological\_transfer\par

\textbf{Dependence of claims.} Regulatory transient claims depend on the selected family/timescale and gain; constant full stimulus equilibrium is independent of tau within R1.\par

\textbf{Discriminating measurement.} Measure this quantity or its identifiable combination in the same preparation and stimulus protocol, with independent WT data; determine physiological ranges before making global robustness claims.\par

\textbf{Source.} \nolinkurl{src/modern_full_model/camp_pka.py:90}. Pena Munzenmayer et al. 2021 gives pathway and endpoint activation evidence; acute kinetic times and exact common capacity gain remain effective settings.\par}

\subsubsection*{Record 92: fully activated increment}
{\small \nolinkurl{regulation.fully_activated_increment}\par

Value: \texttt{0.25}; units: dimensionless; active: Yes; classification: effective\_assumption.\par

Default: ; differs from default: .\par

\textbf{Equation role.} da/dt=(beta-a)/tau; AE4 multiplier=basal+increment*coupling*construct\_scale*a.\par

\textbf{Constraint.} Effective normalisation or gain; pathway activation evidence does not identify each coefficient independently.\par

\textbf{Limit of justification.} An accepted physiological point does not identify this setting or supply a numerical uncertainty interval. No knockout fit or new measurement is performed in this audit.\par

\textbf{Logical domain.} Positive for the retained law; no finite upper physiological bound established.\par

\textbf{Uncertainty status.} genuinely\_unknown\_for\_physiological\_transfer\par

\textbf{Dependence of claims.} Regulatory transient claims depend on the selected family/timescale and gain; constant full stimulus equilibrium is independent of tau within R1. Local numerical dependence is quantified in the linked crosswalk; it establishes neither a measured uncertainty range nor global robustness.\par

\textbf{Discriminating measurement.} Measure the AE4 transport increment under matched stimulation and substrate gradients; resolve the discrepant approximate 2021 gain readings.\par

\textbf{Source.} \nolinkurl{src/modern_full_model/camp_pka.py:90}. Pena Munzenmayer et al. 2021 gives pathway and endpoint activation evidence; acute kinetic times and exact common capacity gain remain effective settings.\par}

\subsubsection*{Record 93: coupling scale}
{\small \nolinkurl{regulation.coupling_scale}\par

Value: \texttt{1}; units: dimensionless; active: Yes; classification: effective\_assumption.\par

Default: ; differs from default: .\par

\textbf{Equation role.} da/dt=(beta-a)/tau; AE4 multiplier=basal+increment*coupling*construct\_scale*a.\par

\textbf{Constraint.} Effective normalisation or gain; pathway activation evidence does not identify each coefficient independently.\par

\textbf{Limit of justification.} An accepted physiological point does not identify this setting or supply a numerical uncertainty interval. No knockout fit or new measurement is performed in this audit.\par

\textbf{Logical domain.} Positive for the retained law; no finite upper physiological bound established.\par

\textbf{Uncertainty status.} genuinely\_unknown\_for\_physiological\_transfer\par

\textbf{Dependence of claims.} Regulatory transient claims depend on the selected family/timescale and gain; constant full stimulus equilibrium is independent of tau within R1.\par

\textbf{Discriminating measurement.} Measure this quantity or its identifiable combination in the same preparation and stimulus protocol, with independent WT data; determine physiological ranges before making global robustness claims.\par

\textbf{Source.} \nolinkurl{src/modern_full_model/camp_pka.py:90}. Pena Munzenmayer et al. 2021 gives pathway and endpoint activation evidence; acute kinetic times and exact common capacity gain remain effective settings.\par}

\subsubsection*{Record 94: fully activated multiplier}
{\small \nolinkurl{nkcc_activation.fully_activated_multiplier}\par

Value: \texttt{1.75}; units: dimensionless; active: Yes; classification: inherited\_effective\_setting.\par

Default: ; differs from default: .\par

\textbf{Equation role.} m\_N=1+(m\_max-1)*clip((Ca-Ca\_rest)/(Ca\_sat-Ca\_rest),0,1).\par

\textbf{Constraint.} Inherited algebraic capacity modulation. Its 0.10 uM saturation endpoint is distinct from the NBC 0.25 uM endpoint; at the production input both are fully recruited.\par

\textbf{Limit of justification.} An accepted physiological point does not identify this setting or supply a numerical uncertainty interval. No knockout fit or new measurement is performed in this audit.\par

\textbf{Logical domain.} Positive for the retained law; no finite upper physiological bound established.\par

\textbf{Uncertainty status.} genuinely\_unknown\_for\_physiological\_transfer\par

\textbf{Dependence of claims.} Numerical state, current and secretion claims are conditional on this active setting; exact stoichiometric identities are independent of its numerical value. Local numerical dependence is quantified in the linked crosswalk; it establishes neither a measured uncertainty range nor global robustness.\par

\textbf{Discriminating measurement.} Measure this quantity or its identifiable combination in the same preparation and stimulus protocol, with independent WT data; determine physiological ranges before making global robustness claims.\par

\textbf{Source.} \nolinkurl{src/modern_full_model/nkcc_stimulation.py:78}. No independently verified primary measurement attached to this numerical setting; model lineage is the source.\par}

\subsubsection*{Record 95: resting calcium uM}
{\small \nolinkurl{nkcc_activation.resting_calcium_uM}\par

Value: \texttt{0.058}; units: uM; active: Yes; classification: inherited\_effective\_setting.\par

Default: ; differs from default: .\par

\textbf{Equation role.} m\_N=1+(m\_max-1)*clip((Ca-Ca\_rest)/(Ca\_sat-Ca\_rest),0,1).\par

\textbf{Constraint.} Inherited algebraic capacity modulation. Its 0.10 uM saturation endpoint is distinct from the NBC 0.25 uM endpoint; at the production input both are fully recruited.\par

\textbf{Limit of justification.} An accepted physiological point does not identify this setting or supply a numerical uncertainty interval. No knockout fit or new measurement is performed in this audit.\par

\textbf{Logical domain.} Positive for the retained law; no finite upper physiological bound established.\par

\textbf{Uncertainty status.} genuinely\_unknown\_for\_physiological\_transfer\par

\textbf{Dependence of claims.} Numerical state, current and secretion claims are conditional on this active setting; exact stoichiometric identities are independent of its numerical value.\par

\textbf{Discriminating measurement.} Measure this quantity or its identifiable combination in the same preparation and stimulus protocol, with independent WT data; determine physiological ranges before making global robustness claims.\par

\textbf{Source.} \nolinkurl{src/modern_full_model/nkcc_stimulation.py:78}. No independently verified primary measurement attached to this numerical setting; model lineage is the source.\par}

\subsubsection*{Record 96: stimulated calcium uM}
{\small \nolinkurl{nkcc_activation.stimulated_calcium_uM}\par

Value: \texttt{0.1}; units: uM; active: Yes; classification: inherited\_effective\_setting.\par

Default: ; differs from default: .\par

\textbf{Equation role.} m\_N=1+(m\_max-1)*clip((Ca-Ca\_rest)/(Ca\_sat-Ca\_rest),0,1).\par

\textbf{Constraint.} Inherited algebraic capacity modulation. Its 0.10 uM saturation endpoint is distinct from the NBC 0.25 uM endpoint; at the production input both are fully recruited.\par

\textbf{Limit of justification.} An accepted physiological point does not identify this setting or supply a numerical uncertainty interval. No knockout fit or new measurement is performed in this audit.\par

\textbf{Logical domain.} Positive for the retained law; no finite upper physiological bound established.\par

\textbf{Uncertainty status.} genuinely\_unknown\_for\_physiological\_transfer\par

\textbf{Dependence of claims.} Numerical state, current and secretion claims are conditional on this active setting; exact stoichiometric identities are independent of its numerical value.\par

\textbf{Discriminating measurement.} Measure this quantity or its identifiable combination in the same preparation and stimulus protocol, with independent WT data; determine physiological ranges before making global robustness claims.\par

\textbf{Source.} \nolinkurl{src/modern_full_model/nkcc_stimulation.py:78}. No independently verified primary measurement attached to this numerical setting; model lineage is the source.\par}

\subsubsection*{Record 97: arm}
{\small \nolinkurl{protocol.arm}\par

Value: CCH\_IPR; units: configuration; active: Yes; classification: protocol\_setting.\par

Default: ; differs from default: .\par

\textbf{Equation role.} Declared input arm and time window select calcium and beta activation steps; no agonist dose response map is present.\par

\textbf{Constraint.} Declared experimental-input representation, including a step waveform; this is not evidence for measured intracellular calcium kinetics.\par

\textbf{Limit of justification.} An accepted physiological point does not identify this setting or supply a numerical uncertainty interval. No knockout fit or new measurement is performed in this audit.\par

\textbf{Logical domain.} Selected discrete setting or stated numerical reference.\par

\textbf{Uncertainty status.} genuinely\_unknown\_for\_physiological\_transfer\par

\textbf{Dependence of claims.} Numerical state, current and secretion claims are conditional on this active setting; exact stoichiometric identities are independent of its numerical value.\par

\textbf{Discriminating measurement.} Measure this quantity or its identifiable combination in the same preparation and stimulus protocol, with independent WT data; determine physiological ranges before making global robustness claims.\par

\textbf{Source.} \nolinkurl{src/modern_full_model/validation.py:135}. No independently verified primary measurement attached to this numerical setting; model lineage is the source.\par}

\subsubsection*{Record 98: onset s}
{\small \nolinkurl{protocol.onset_s}\par

Value: \texttt{0}; units: s; active: Yes; classification: protocol\_setting.\par

Default: ; differs from default: .\par

\textbf{Equation role.} Declared input arm and time window select calcium and beta activation steps; no agonist dose response map is present.\par

\textbf{Constraint.} Declared experimental-input representation, including a step waveform; this is not evidence for measured intracellular calcium kinetics.\par

\textbf{Limit of justification.} An accepted physiological point does not identify this setting or supply a numerical uncertainty interval. No knockout fit or new measurement is performed in this audit.\par

\textbf{Logical domain.} Exact zero records a declared omitted contribution or protocol normalisation; it is not a confidence bound.\par

\textbf{Uncertainty status.} genuinely\_unknown\_for\_physiological\_transfer\par

\textbf{Dependence of claims.} Numerical state, current and secretion claims are conditional on this active setting; exact stoichiometric identities are independent of its numerical value.\par

\textbf{Discriminating measurement.} Measure this quantity or its identifiable combination in the same preparation and stimulus protocol, with independent WT data; determine physiological ranges before making global robustness claims.\par

\textbf{Source.} \nolinkurl{src/modern_full_model/validation.py:135}. No independently verified primary measurement attached to this numerical setting; model lineage is the source.\par}

\subsubsection*{Record 99: duration s}
{\small \nolinkurl{protocol.duration_s}\par

Value: \texttt{600}; units: s; active: Yes; classification: protocol\_setting.\par

Default: ; differs from default: .\par

\textbf{Equation role.} Declared input arm and time window select calcium and beta activation steps; no agonist dose response map is present.\par

\textbf{Constraint.} Declared experimental-input representation, including a step waveform; this is not evidence for measured intracellular calcium kinetics.\par

\textbf{Limit of justification.} An accepted physiological point does not identify this setting or supply a numerical uncertainty interval. No knockout fit or new measurement is performed in this audit.\par

\textbf{Logical domain.} Positive for the retained law; no finite upper physiological bound established.\par

\textbf{Uncertainty status.} genuinely\_unknown\_for\_physiological\_transfer\par

\textbf{Dependence of claims.} Numerical state, current and secretion claims are conditional on this active setting; exact stoichiometric identities are independent of its numerical value.\par

\textbf{Discriminating measurement.} Measure this quantity or its identifiable combination in the same preparation and stimulus protocol, with independent WT data; determine physiological ranges before making global robustness claims.\par

\textbf{Source.} \nolinkurl{src/modern_full_model/validation.py:135}. 2015 protocol lineage recorded by WTProtocol.provenance in validation.py; implementation transfer retains dose labels and duration.\par}

\subsubsection*{Record 100: cch dose uM}
{\small \nolinkurl{protocol.cch_dose_uM}\par

Value: \texttt{0.3}; units: uM; active: No; classification: protocol\_setting.\par

Default: ; differs from default: .\par

\textbf{Equation role.} Declared input arm and time window select calcium and beta activation steps; no agonist dose response map is present.\par

\textbf{Constraint.} Nominal experiment dose labels retained for provenance. The selected model has no dose response map from these fields to calcium or beta activation.\par

\textbf{Limit of justification.} An accepted physiological point does not identify this setting or supply a numerical uncertainty interval. No knockout fit or new measurement is performed in this audit.\par

\textbf{Logical domain.} Positive for the retained law; no finite upper physiological bound established.\par

\textbf{Uncertainty status.} genuinely\_unknown\_for\_physiological\_transfer\par

\textbf{Dependence of claims.} No independent parent RHS sensitivity. Used only for the explicitly stated metadata, initial reference, design lineage or inverse family role.\par

\textbf{Discriminating measurement.} Measure this quantity or its identifiable combination in the same preparation and stimulus protocol, with independent WT data; determine physiological ranges before making global robustness claims.\par

\textbf{Source.} \nolinkurl{src/modern_full_model/validation.py:135}. 2015 protocol lineage recorded by WTProtocol.provenance in validation.py; implementation transfer retains dose labels and duration.\par}

\subsubsection*{Record 101: ipr dose uM}
{\small \nolinkurl{protocol.ipr_dose_uM}\par

Value: \texttt{5}; units: uM; active: No; classification: protocol\_setting.\par

Default: ; differs from default: .\par

\textbf{Equation role.} Declared input arm and time window select calcium and beta activation steps; no agonist dose response map is present.\par

\textbf{Constraint.} Nominal experiment dose labels retained for provenance. The selected model has no dose response map from these fields to calcium or beta activation.\par

\textbf{Limit of justification.} An accepted physiological point does not identify this setting or supply a numerical uncertainty interval. No knockout fit or new measurement is performed in this audit.\par

\textbf{Logical domain.} Positive for the retained law; no finite upper physiological bound established.\par

\textbf{Uncertainty status.} genuinely\_unknown\_for\_physiological\_transfer\par

\textbf{Dependence of claims.} No independent parent RHS sensitivity. Used only for the explicitly stated metadata, initial reference, design lineage or inverse family role.\par

\textbf{Discriminating measurement.} Measure this quantity or its identifiable combination in the same preparation and stimulus protocol, with independent WT data; determine physiological ranges before making global robustness claims.\par

\textbf{Source.} \nolinkurl{src/modern_full_model/validation.py:135}. 2015 protocol lineage recorded by WTProtocol.provenance in validation.py; implementation transfer retains dose labels and duration.\par}

\subsubsection*{Record 102: resting calcium uM}
{\small \nolinkurl{protocol.resting_calcium_uM}\par

Value: \texttt{0.058}; units: uM; active: Yes; classification: protocol\_setting.\par

Default: ; differs from default: .\par

\textbf{Equation role.} Declared input arm and time window select calcium and beta activation steps; no agonist dose response map is present.\par

\textbf{Constraint.} Declared experimental-input representation, including a step waveform; this is not evidence for measured intracellular calcium kinetics.\par

\textbf{Limit of justification.} An accepted physiological point does not identify this setting or supply a numerical uncertainty interval. No knockout fit or new measurement is performed in this audit.\par

\textbf{Logical domain.} Positive for the retained law; no finite upper physiological bound established.\par

\textbf{Uncertainty status.} genuinely\_unknown\_for\_physiological\_transfer\par

\textbf{Dependence of claims.} Numerical state, current and secretion claims are conditional on this active setting; exact stoichiometric identities are independent of its numerical value.\par

\textbf{Discriminating measurement.} Measure this quantity or its identifiable combination in the same preparation and stimulus protocol, with independent WT data; determine physiological ranges before making global robustness claims.\par

\textbf{Source.} \nolinkurl{src/modern_full_model/validation.py:135}. No independently verified primary measurement attached to this numerical setting; model lineage is the source.\par}

\subsubsection*{Record 103: stimulated calcium uM}
{\small \nolinkurl{protocol.stimulated_calcium_uM}\par

Value: \texttt{0.25}; units: uM; active: Yes; classification: protocol\_setting.\par

Default: ; differs from default: .\par

\textbf{Equation role.} Declared input arm and time window select calcium and beta activation steps; no agonist dose response map is present.\par

\textbf{Constraint.} Declared experimental-input representation, including a step waveform; this is not evidence for measured intracellular calcium kinetics.\par

\textbf{Limit of justification.} An accepted physiological point does not identify this setting or supply a numerical uncertainty interval. No knockout fit or new measurement is performed in this audit.\par

\textbf{Logical domain.} Positive for the retained law; no finite upper physiological bound established.\par

\textbf{Uncertainty status.} genuinely\_unknown\_for\_physiological\_transfer\par

\textbf{Dependence of claims.} Numerical state, current and secretion claims are conditional on this active setting; exact stoichiometric identities are independent of its numerical value.\par

\textbf{Discriminating measurement.} Measure this quantity or its identifiable combination in the same preparation and stimulus protocol, with independent WT data; determine physiological ranges before making global robustness claims.\par

\textbf{Source.} \nolinkurl{src/modern_full_model/validation.py:135}. No independently verified primary measurement attached to this numerical setting; model lineage is the source.\par}

\subsubsection*{Record 104: beta occupancy on}
{\small \nolinkurl{protocol.beta_occupancy_on}\par

Value: \texttt{1}; units: dimensionless fraction; active: Yes; classification: protocol\_setting.\par

Default: ; differs from default: .\par

\textbf{Equation role.} Declared input arm and time window select calcium and beta activation steps; no agonist dose response map is present.\par

\textbf{Constraint.} Declared experimental-input representation, including a step waveform; this is not evidence for measured intracellular calcium kinetics.\par

\textbf{Limit of justification.} An accepted physiological point does not identify this setting or supply a numerical uncertainty interval. No knockout fit or new measurement is performed in this audit.\par

\textbf{Logical domain.} [0,1] by fractional definition only; not a measured uncertainty interval.\par

\textbf{Uncertainty status.} genuinely\_unknown\_for\_physiological\_transfer\par

\textbf{Dependence of claims.} Numerical state, current and secretion claims are conditional on this active setting; exact stoichiometric identities are independent of its numerical value.\par

\textbf{Discriminating measurement.} Measure this quantity or its identifiable combination in the same preparation and stimulus protocol, with independent WT data; determine physiological ranges before making global robustness claims.\par

\textbf{Source.} \nolinkurl{src/modern_full_model/validation.py:135}. No independently verified primary measurement attached to this numerical setting; model lineage is the source.\par}

\subsubsection*{Record 105: b}
{\small \nolinkurl{Task41.b}\par

Value: \texttt{0.1051187284328845}; units: dimensionless; active: No; classification: target\_selected.\par

Default: ; differs from default: .\par

\textbf{Equation role.} In the prescribed inverse family, residual stimulated CaCC recruitment is b and the AE4 dependent fraction is rho=1-b.\par

\textbf{Constraint.} Residual CaCC recruitment in the deliberately inverse Task 41 construction. Not part of the parent Task 40 parameter set.\par

\textbf{Limit of justification.} Matching the target secretion comparison is inverse selection, not independent validation or experimental estimation of CaCC recruitment.\par

\textbf{Logical domain.} 0 < b <= 1, equivalently 0 <= rho < 1, by the prescribed family constructor. b=0 is excluded; this is not a measured uncertainty interval.\par

\textbf{Uncertainty status.} genuinely\_unknown\_for\_physiological\_transfer\par

\textbf{Dependence of claims.} Directly controls the target selected Task 41 CaCC recruitment mechanism and its 600 s inverse crossing; outside the parent Task 40 model.\par

\textbf{Discriminating measurement.} Measure the fraction of stimulated CaCC recruitment lost with acute AE4 depletion while controlling calcium, pH, chloride gradient, voltage and expression; test the prescribed family rather than use the secretion target again.\par

\textbf{Source.} \nolinkurl{src/modern_full_model/ae4_cacc_recruitment.py:24}. No independently verified primary measurement attached to this numerical setting; model lineage is the source.\par}

\subsubsection*{Record 106: intracellular modifier hill}
{\small \nolinkurl{cha.intracellular_modifier_hill}\par

Value: \texttt{3}; units: dimensionless; active: Yes; classification: published\_kinetic\_coefficient.\par

Default: ; differs from default: .\par

\textbf{Equation role.} JH=N\_H*m\_H*Mod2*(a*b-c*d)/(a+b+c+d), with rate conversion from ms to s and printed fourth reverse rate retained.\par

\textbf{Constraint.} Cha 2009 Eq. 5 Mod2 exponent; a fixed published exponent, not an identified salivary abundance.\par

\textbf{Limit of justification.} Publication of the kinetic law does not measure salivary carrier abundance or establish an uncertainty interval for transfer to this preparation.\par

\textbf{Logical domain.} Positive for the retained law; no finite upper physiological bound established.\par

\textbf{Uncertainty status.} genuinely\_unknown\_for\_physiological\_transfer\par

\textbf{Dependence of claims.} Numerical state, current and secretion claims are conditional on this active setting; exact stoichiometric identities are independent of its numerical value.\par

\textbf{Discriminating measurement.} Measure this quantity or its identifiable combination in the same preparation and stimulus protocol, with independent WT data; determine physiological ranges before making global robustness claims.\par

\textbf{Source.} \nolinkurl{src/modern_full_model/nhe1_cha2009.py:113}. Cha et al. 2009, Eqs. 3 and 5, Table S1; DOI 10.1016/j.bpj.2009.08.053; repository results/31\_nhe1\_mechanistic\_repair/source\_table\_s1.json.\par}

\subsubsection*{Record 107: extracellular modifier hill}
{\small \nolinkurl{cha.extracellular_modifier_hill}\par

Value: \texttt{1}; units: dimensionless; active: Yes; classification: published\_kinetic\_coefficient.\par

Default: ; differs from default: .\par

\textbf{Equation role.} JH=N\_H*m\_H*Mod2*(a*b-c*d)/(a+b+c+d), with rate conversion from ms to s and printed fourth reverse rate retained.\par

\textbf{Constraint.} Cha 2009 Eq. 5 Mod2 exponent; a fixed published exponent, not an identified salivary abundance.\par

\textbf{Limit of justification.} Publication of the kinetic law does not measure salivary carrier abundance or establish an uncertainty interval for transfer to this preparation.\par

\textbf{Logical domain.} Positive for the retained law; no finite upper physiological bound established.\par

\textbf{Uncertainty status.} genuinely\_unknown\_for\_physiological\_transfer\par

\textbf{Dependence of claims.} Numerical state, current and secretion claims are conditional on this active setting; exact stoichiometric identities are independent of its numerical value.\par

\textbf{Discriminating measurement.} Measure this quantity or its identifiable combination in the same preparation and stimulus protocol, with independent WT data; determine physiological ranges before making global robustness claims.\par

\textbf{Source.} \nolinkurl{src/modern_full_model/nhe1_cha2009.py:113}. Cha et al. 2009, Eqs. 3 and 5, Table S1; DOI 10.1016/j.bpj.2009.08.053; repository results/31\_nhe1\_mechanistic\_repair/source\_table\_s1.json.\par}

\subsubsection*{Record 108: law}
{\small \nolinkurl{nkcc.law}\par

Value: palk\_benjamin\_2010\_eq17; units: model identifier; active: Yes; classification: law\_selection.\par

Default: ; differs from default: .\par

\textbf{Equation role.} JN=alpha*m\_N*(A1-A2[Na]i[K]i[Cl]i\textasciicircum{}2)/(A3+A4[Na]i[K]i[Cl]i\textasciicircum{}2); inward 1 Na, 1 K, 2 Cl per cycle.\par

\textbf{Constraint.} Explicit Palk/Benjamin law selected by the Task 39 builder and inherited by Task 40.\par

\textbf{Limit of justification.} An accepted physiological point does not identify this setting or supply a numerical uncertainty interval. No knockout fit or new measurement is performed in this audit.\par

\textbf{Logical domain.} Selected discrete setting or stated numerical reference.\par

\textbf{Uncertainty status.} not\_a\_measured\_uncertainty\_parameter\par

\textbf{Dependence of claims.} Numerical state, current and secretion claims are conditional on this active setting; exact stoichiometric identities are independent of its numerical value.\par

\textbf{Discriminating measurement.} Measure this quantity or its identifiable combination in the same preparation and stimulus protocol, with independent WT data; determine physiological ranges before making global robustness claims.\par

\textbf{Source.} \nolinkurl{src/modern_full_model/nkcc1_palk2010.py:31}. Palk et al. 2010/Benjamin and Johnson 1997 law, reproduced as Vera Siguenza et al. 2018 Eq. 17; coefficients in src/modern\_full\_model/nkcc1\_palk2010.py.\par}

\subsubsection*{Record 109: law}
{\small \nolinkurl{ae4.law}\par

Value: equal\_routing\_adapter; units: model identifier; active: Yes; classification: law\_selection.\par

Default: ; differs from default: .\par

\textbf{Equation role.} Shared carrier QSS determines scalar J4; equal routing imposes (-J4/2,-J4/2,J4,-2J4,-2J4) in cell Na,K,Cl,TIC,TA.\par

\textbf{Constraint.} Equal cation routing fixes source shares at one half each while retaining the inherited total scalar cycle flux.\par

\textbf{Limit of justification.} An accepted physiological point does not identify this setting or supply a numerical uncertainty interval. No knockout fit or new measurement is performed in this audit.\par

\textbf{Logical domain.} Selected discrete setting or stated numerical reference.\par

\textbf{Uncertainty status.} not\_a\_measured\_uncertainty\_parameter\par

\textbf{Dependence of claims.} Numerical state, current and secretion claims are conditional on this active setting; exact stoichiometric identities are independent of its numerical value.\par

\textbf{Discriminating measurement.} Measure this quantity or its identifiable combination in the same preparation and stimulus protocol, with independent WT data; determine physiological ranges before making global robustness claims.\par

\textbf{Source.} \nolinkurl{src/modern_full_model/transporters.py:458}. No independently verified primary measurement attached to this numerical setting; model lineage is the source.\par}

\subsubsection*{Record 110: family}
{\small \nolinkurl{regulation.family}\par

Value: R1; units: model identifier; active: Yes; classification: law\_selection.\par

Default: ; differs from default: .\par

\textbf{Equation role.} da/dt=(beta-a)/tau; AE4 multiplier=basal+increment*coupling*construct\_scale*a.\par

\textbf{Constraint.} One retained effective AE4 activation state; no measured separation of cAMP and PKA kinetics.\par

\textbf{Limit of justification.} An accepted physiological point does not identify this setting or supply a numerical uncertainty interval. No knockout fit or new measurement is performed in this audit.\par

\textbf{Logical domain.} Selected discrete setting or stated numerical reference.\par

\textbf{Uncertainty status.} not\_a\_measured\_uncertainty\_parameter\par

\textbf{Dependence of claims.} Regulatory transient claims depend on the selected family/timescale and gain; constant full stimulus equilibrium is independent of tau within R1.\par

\textbf{Discriminating measurement.} Measure this quantity or its identifiable combination in the same preparation and stimulus protocol, with independent WT data; determine physiological ranges before making global robustness claims.\par

\textbf{Source.} \nolinkurl{src/modern_full_model/camp_pka.py:90}. Pena Munzenmayer et al. 2021 gives pathway and endpoint activation evidence; acute kinetic times and exact common capacity gain remain effective settings.\par}

\subsubsection*{Record 111: construct}
{\small \nolinkurl{regulation.construct}\par

Value: WT; units: construct identifier; active: Yes; classification: protocol\_setting.\par

Default: ; differs from default: .\par

\textbf{Equation role.} da/dt=(beta-a)/tau; AE4 multiplier=basal+increment*coupling*construct\_scale*a.\par

\textbf{Constraint.} WT construct selects unit response scale in the gain map.\par

\textbf{Limit of justification.} An accepted physiological point does not identify this setting or supply a numerical uncertainty interval. No knockout fit or new measurement is performed in this audit.\par

\textbf{Logical domain.} Selected discrete setting or stated numerical reference.\par

\textbf{Uncertainty status.} genuinely\_unknown\_for\_physiological\_transfer\par

\textbf{Dependence of claims.} Regulatory transient claims depend on the selected family/timescale and gain; constant full stimulus equilibrium is independent of tau within R1.\par

\textbf{Discriminating measurement.} Measure this quantity or its identifiable combination in the same preparation and stimulus protocol, with independent WT data; determine physiological ranges before making global robustness claims.\par

\textbf{Source.} \nolinkurl{src/modern_full_model/camp_pka.py:90}. Pena Munzenmayer et al. 2021 gives pathway and endpoint activation evidence; acute kinetic times and exact common capacity gain remain effective settings.\par}

\subsubsection*{Record 112: family}
{\small \nolinkurl{nkcc_activation.family}\par

Value: N1\_CCH\_ALGEBRAIC; units: model identifier; active: Yes; classification: law\_selection.\par

Default: ; differs from default: .\par

\textbf{Equation role.} m\_N=1+(m\_max-1)*clip((Ca-Ca\_rest)/(Ca\_sat-Ca\_rest),0,1).\par

\textbf{Constraint.} Selected instantaneous calcium recruitment family; no added NKCC dynamic state.\par

\textbf{Limit of justification.} An accepted physiological point does not identify this setting or supply a numerical uncertainty interval. No knockout fit or new measurement is performed in this audit.\par

\textbf{Logical domain.} Selected discrete setting or stated numerical reference.\par

\textbf{Uncertainty status.} not\_a\_measured\_uncertainty\_parameter\par

\textbf{Dependence of claims.} Numerical state, current and secretion claims are conditional on this active setting; exact stoichiometric identities are independent of its numerical value.\par

\textbf{Discriminating measurement.} Measure this quantity or its identifiable combination in the same preparation and stimulus protocol, with independent WT data; determine physiological ranges before making global robustness claims.\par

\textbf{Source.} \nolinkurl{src/modern_full_model/nkcc_stimulation.py:78}. No independently verified primary measurement attached to this numerical setting; model lineage is the source.\par}

\subsubsection*{Record 113: TARGET NKCC1 POSITIVE CL SHARE}
{\small \nolinkurl{nbc_design.TARGET_NKCC1_POSITIVE_CL_SHARE}\par

Value: \texttt{0.7}; units: dimensionless; active: No; classification: effective\_assumption.\par

Default: ; differs from default: .\par

\textbf{Equation role.} Upstream construction: J4\_req=((1-s)/s)*L\_N; JB\_req=J4\_req-H\_stim/2; capacity includes the self consistent voltage.\par

\textbf{Constraint.} Design input or conditional derived quantity upstream of the frozen NBC capacity; not independently evaluated in the production RHS. The 70/30 share is assumed, and the reference fluxes are modelled.\par

\textbf{Limit of justification.} The 70/30 partition is assumed, not measured. Solving the coupled current balance does not establish the physiological truth or unique identification of that partition.\par

\textbf{Logical domain.} [0,1] by fractional definition only; not a measured uncertainty interval.\par

\textbf{Uncertainty status.} conditional\_design\_value\_no\_experimental\_interval\par

\textbf{Dependence of claims.} No independent parent RHS sensitivity. Used only for the explicitly stated metadata, initial reference, design lineage or inverse family role.\par

\textbf{Discriminating measurement.} Measure this quantity or its identifiable combination in the same preparation and stimulus protocol, with independent WT data; determine physiological ranges before making global robustness claims.\par

\textbf{Source.} \nolinkurl{src/modern_full_model/nbc_minimal.py:77}. No independently verified primary measurement attached to this numerical setting; model lineage is the source.\par}

\subsubsection*{Record 114: TASK31 R09 NKCC1 CL FMOL S}
{\small \nolinkurl{nbc_design.TASK31_R09_NKCC1_CL_FMOL_S}\par

Value: \texttt{0.2562444047732706}; units: fmol/s; active: No; classification: WT\_architecture\_derived.\par

Default: ; differs from default: .\par

\textbf{Equation role.} Upstream construction: J4\_req=((1-s)/s)*L\_N; JB\_req=J4\_req-H\_stim/2; capacity includes the self consistent voltage.\par

\textbf{Constraint.} Design input or conditional derived quantity upstream of the frozen NBC capacity; not independently evaluated in the production RHS. The 70/30 share is assumed, and the reference fluxes are modelled.\par

\textbf{Limit of justification.} The 70/30 partition is assumed, not measured. Solving the coupled current balance does not establish the physiological truth or unique identification of that partition.\par

\textbf{Logical domain.} Positive for the retained law; no finite upper physiological bound established.\par

\textbf{Uncertainty status.} conditional\_design\_value\_no\_experimental\_interval\par

\textbf{Dependence of claims.} No independent parent RHS sensitivity. Used only for the explicitly stated metadata, initial reference, design lineage or inverse family role.\par

\textbf{Discriminating measurement.} Measure this quantity or its identifiable combination in the same preparation and stimulus protocol, with independent WT data; determine physiological ranges before making global robustness claims.\par

\textbf{Source.} \nolinkurl{src/modern_full_model/nbc_minimal.py:77}. No independently verified primary measurement attached to this numerical setting; model lineage is the source.\par}

\subsubsection*{Record 115: TASK31 R09 NHE1 FMOL S}
{\small \nolinkurl{nbc_design.TASK31_R09_NHE1_FMOL_S}\par

Value: \texttt{0.007890580489555009}; units: fmol/s; active: No; classification: WT\_architecture\_derived.\par

Default: ; differs from default: .\par

\textbf{Equation role.} Upstream construction: J4\_req=((1-s)/s)*L\_N; JB\_req=J4\_req-H\_stim/2; capacity includes the self consistent voltage.\par

\textbf{Constraint.} Design input or conditional derived quantity upstream of the frozen NBC capacity; not independently evaluated in the production RHS. The 70/30 share is assumed, and the reference fluxes are modelled.\par

\textbf{Limit of justification.} The 70/30 partition is assumed, not measured. Solving the coupled current balance does not establish the physiological truth or unique identification of that partition.\par

\textbf{Logical domain.} Positive for the retained law; no finite upper physiological bound established.\par

\textbf{Uncertainty status.} conditional\_design\_value\_no\_experimental\_interval\par

\textbf{Dependence of claims.} No independent parent RHS sensitivity. Used only for the explicitly stated metadata, initial reference, design lineage or inverse family role.\par

\textbf{Discriminating measurement.} Measure this quantity or its identifiable combination in the same preparation and stimulus protocol, with independent WT data; determine physiological ranges before making global robustness claims.\par

\textbf{Source.} \nolinkurl{src/modern_full_model/nbc_minimal.py:77}. No independently verified primary measurement attached to this numerical setting; model lineage is the source.\par}

\subsubsection*{Record 116: DERIVED AE4 CL FMOL S}
{\small \nolinkurl{nbc_design.DERIVED_AE4_CL_FMOL_S}\par

Value: \texttt{0.109819030617116}; units: fmol/s; active: No; classification: WT\_architecture\_derived.\par

Default: ; differs from default: .\par

\textbf{Equation role.} Upstream construction: J4\_req=((1-s)/s)*L\_N; JB\_req=J4\_req-H\_stim/2; capacity includes the self consistent voltage.\par

\textbf{Constraint.} Design input or conditional derived quantity upstream of the frozen NBC capacity; not independently evaluated in the production RHS. The 70/30 share is assumed, and the reference fluxes are modelled.\par

\textbf{Limit of justification.} The 70/30 partition is assumed, not measured. Solving the coupled current balance does not establish the physiological truth or unique identification of that partition.\par

\textbf{Logical domain.} Positive for the retained law; no finite upper physiological bound established.\par

\textbf{Uncertainty status.} conditional\_design\_value\_no\_experimental\_interval\par

\textbf{Dependence of claims.} No independent parent RHS sensitivity. Used only for the explicitly stated metadata, initial reference, design lineage or inverse family role.\par

\textbf{Discriminating measurement.} Measure this quantity or its identifiable combination in the same preparation and stimulus protocol, with independent WT data; determine physiological ranges before making global robustness claims.\par

\textbf{Source.} \nolinkurl{src/modern_full_model/nbc_minimal.py:77}. No independently verified primary measurement attached to this numerical setting; model lineage is the source.\par}

\subsubsection*{Record 117: DERIVED STIMULATED NHE1 FMOL S}
{\small \nolinkurl{nbc_design.DERIVED_STIMULATED_NHE1_FMOL_S}\par

Value: \texttt{0.01814833512597652}; units: fmol/s; active: No; classification: WT\_architecture\_derived.\par

Default: ; differs from default: .\par

\textbf{Equation role.} Upstream construction: J4\_req=((1-s)/s)*L\_N; JB\_req=J4\_req-H\_stim/2; capacity includes the self consistent voltage.\par

\textbf{Constraint.} Design input or conditional derived quantity upstream of the frozen NBC capacity; not independently evaluated in the production RHS. The 70/30 share is assumed, and the reference fluxes are modelled.\par

\textbf{Limit of justification.} The 70/30 partition is assumed, not measured. Solving the coupled current balance does not establish the physiological truth or unique identification of that partition.\par

\textbf{Logical domain.} Positive for the retained law; no finite upper physiological bound established.\par

\textbf{Uncertainty status.} conditional\_design\_value\_no\_experimental\_interval\par

\textbf{Dependence of claims.} No independent parent RHS sensitivity. Used only for the explicitly stated metadata, initial reference, design lineage or inverse family role.\par

\textbf{Discriminating measurement.} Measure this quantity or its identifiable combination in the same preparation and stimulus protocol, with independent WT data; determine physiological ranges before making global robustness claims.\par

\textbf{Source.} \nolinkurl{src/modern_full_model/nbc_minimal.py:77}. No independently verified primary measurement attached to this numerical setting; model lineage is the source.\par}

\subsubsection*{Record 118: DERIVED REQUIRED NBC CYCLE FMOL S}
{\small \nolinkurl{nbc_design.DERIVED_REQUIRED_NBC_CYCLE_FMOL_S}\par

Value: \texttt{0.1007448630541277}; units: fmol/s; active: No; classification: WT\_architecture\_derived.\par

Default: ; differs from default: .\par

\textbf{Equation role.} Upstream construction: J4\_req=((1-s)/s)*L\_N; JB\_req=J4\_req-H\_stim/2; capacity includes the self consistent voltage.\par

\textbf{Constraint.} Design input or conditional derived quantity upstream of the frozen NBC capacity; not independently evaluated in the production RHS. The 70/30 share is assumed, and the reference fluxes are modelled.\par

\textbf{Limit of justification.} The 70/30 partition is assumed, not measured. Solving the coupled current balance does not establish the physiological truth or unique identification of that partition.\par

\textbf{Logical domain.} Positive for the retained law; no finite upper physiological bound established.\par

\textbf{Uncertainty status.} conditional\_design\_value\_no\_experimental\_interval\par

\textbf{Dependence of claims.} No independent parent RHS sensitivity. Used only for the explicitly stated metadata, initial reference, design lineage or inverse family role.\par

\textbf{Discriminating measurement.} Measure this quantity or its identifiable combination in the same preparation and stimulus protocol, with independent WT data; determine physiological ranges before making global robustness claims.\par

\textbf{Source.} \nolinkurl{src/modern_full_model/nbc_minimal.py:77}. No independently verified primary measurement attached to this numerical setting; model lineage is the source.\par}

\subsubsection*{Record 119: na i fmol}
{\small \nolinkurl{frozen_initial.na_i_fmol}\par

Value: \texttt{16.72035873274549}; units: fmol; active: Yes; classification: WT\_constraint\_derived.\par

Default: ; differs from default: .\par

\textbf{Equation role.} Full conserved state and regulatory initial value at t=0 for the finite duration protocol; all amounts and both volumes evolve subsequently.\par

\textbf{Constraint.} Solved, accepted WT resting state used as the shared initial condition for production genotype trajectories; not the constructor seed and not an independent measurement.\par

\textbf{Limit of justification.} An accepted physiological point does not identify this setting or supply a numerical uncertainty interval. No knockout fit or new measurement is performed in this audit.\par

\textbf{Logical domain.} Positive for the retained law; no finite upper physiological bound established.\par

\textbf{Uncertainty status.} conditional\_root\_value\_no\_experimental\_interval\par

\textbf{Dependence of claims.} Numerical state, current and secretion claims are conditional on this active setting; exact stoichiometric identities are independent of its numerical value.\par

\textbf{Discriminating measurement.} Measure this quantity or its identifiable combination in the same preparation and stimulus protocol, with independent WT data; determine physiological ranges before making global robustness claims.\par

\textbf{Source.} \nolinkurl{src/modern_full_model/model.py:559}. No independently verified primary measurement attached to this numerical setting; model lineage is the source.\par}

\subsubsection*{Record 120: k i fmol}
{\small \nolinkurl{frozen_initial.k_i_fmol}\par

Value: \texttt{170.0020085218768}; units: fmol; active: Yes; classification: WT\_constraint\_derived.\par

Default: ; differs from default: .\par

\textbf{Equation role.} Full conserved state and regulatory initial value at t=0 for the finite duration protocol; all amounts and both volumes evolve subsequently.\par

\textbf{Constraint.} Solved, accepted WT resting state used as the shared initial condition for production genotype trajectories; not the constructor seed and not an independent measurement.\par

\textbf{Limit of justification.} An accepted physiological point does not identify this setting or supply a numerical uncertainty interval. No knockout fit or new measurement is performed in this audit.\par

\textbf{Logical domain.} Positive for the retained law; no finite upper physiological bound established.\par

\textbf{Uncertainty status.} conditional\_root\_value\_no\_experimental\_interval\par

\textbf{Dependence of claims.} Numerical state, current and secretion claims are conditional on this active setting; exact stoichiometric identities are independent of its numerical value.\par

\textbf{Discriminating measurement.} Measure this quantity or its identifiable combination in the same preparation and stimulus protocol, with independent WT data; determine physiological ranges before making global robustness claims.\par

\textbf{Source.} \nolinkurl{src/modern_full_model/model.py:559}. No independently verified primary measurement attached to this numerical setting; model lineage is the source.\par}

\subsubsection*{Record 121: cl i fmol}
{\small \nolinkurl{frozen_initial.cl_i_fmol}\par

Value: \texttt{88.03180196356951}; units: fmol; active: Yes; classification: WT\_constraint\_derived.\par

Default: ; differs from default: .\par

\textbf{Equation role.} Full conserved state and regulatory initial value at t=0 for the finite duration protocol; all amounts and both volumes evolve subsequently.\par

\textbf{Constraint.} Solved, accepted WT resting state used as the shared initial condition for production genotype trajectories; not the constructor seed and not an independent measurement.\par

\textbf{Limit of justification.} An accepted physiological point does not identify this setting or supply a numerical uncertainty interval. No knockout fit or new measurement is performed in this audit.\par

\textbf{Logical domain.} Positive for the retained law; no finite upper physiological bound established.\par

\textbf{Uncertainty status.} conditional\_root\_value\_no\_experimental\_interval\par

\textbf{Dependence of claims.} Numerical state, current and secretion claims are conditional on this active setting; exact stoichiometric identities are independent of its numerical value.\par

\textbf{Discriminating measurement.} Measure this quantity or its identifiable combination in the same preparation and stimulus protocol, with independent WT data; determine physiological ranges before making global robustness claims.\par

\textbf{Source.} \nolinkurl{src/modern_full_model/model.py:559}. No independently verified primary measurement attached to this numerical setting; model lineage is the source.\par}

\subsubsection*{Record 122: tic i fmol}
{\small \nolinkurl{frozen_initial.tic_i_fmol}\par

Value: \texttt{8.263511808402752}; units: fmol; active: Yes; classification: WT\_constraint\_derived.\par

Default: ; differs from default: .\par

\textbf{Equation role.} Full conserved state and regulatory initial value at t=0 for the finite duration protocol; all amounts and both volumes evolve subsequently.\par

\textbf{Constraint.} Solved, accepted WT resting state used as the shared initial condition for production genotype trajectories; not the constructor seed and not an independent measurement.\par

\textbf{Limit of justification.} An accepted physiological point does not identify this setting or supply a numerical uncertainty interval. No knockout fit or new measurement is performed in this audit.\par

\textbf{Logical domain.} Positive for the retained law; no finite upper physiological bound established.\par

\textbf{Uncertainty status.} conditional\_root\_value\_no\_experimental\_interval\par

\textbf{Dependence of claims.} Numerical state, current and secretion claims are conditional on this active setting; exact stoichiometric identities are independent of its numerical value.\par

\textbf{Discriminating measurement.} Measure this quantity or its identifiable combination in the same preparation and stimulus protocol, with independent WT data; determine physiological ranges before making global robustness claims.\par

\textbf{Source.} \nolinkurl{src/modern_full_model/model.py:559}. No independently verified primary measurement attached to this numerical setting; model lineage is the source.\par}

\subsubsection*{Record 123: alk i fmol}
{\small \nolinkurl{frozen_initial.alk_i_fmol}\par

Value: \texttt{20.62880308510795}; units: fmol; active: Yes; classification: WT\_constraint\_derived.\par

Default: ; differs from default: .\par

\textbf{Equation role.} Full conserved state and regulatory initial value at t=0 for the finite duration protocol; all amounts and both volumes evolve subsequently.\par

\textbf{Constraint.} Solved, accepted WT resting state used as the shared initial condition for production genotype trajectories; not the constructor seed and not an independent measurement.\par

\textbf{Limit of justification.} An accepted physiological point does not identify this setting or supply a numerical uncertainty interval. No knockout fit or new measurement is performed in this audit.\par

\textbf{Logical domain.} Positive for the retained law; no finite upper physiological bound established.\par

\textbf{Uncertainty status.} conditional\_root\_value\_no\_experimental\_interval\par

\textbf{Dependence of claims.} Numerical state, current and secretion claims are conditional on this active setting; exact stoichiometric identities are independent of its numerical value.\par

\textbf{Discriminating measurement.} Measure this quantity or its identifiable combination in the same preparation and stimulus protocol, with independent WT data; determine physiological ranges before making global robustness claims.\par

\textbf{Source.} \nolinkurl{src/modern_full_model/model.py:559}. No independently verified primary measurement attached to this numerical setting; model lineage is the source.\par}

\subsubsection*{Record 124: volume i pL}
{\small \nolinkurl{frozen_initial.volume_i_pL}\par

Value: \texttt{1.45333439258746}; units: pL; active: Yes; classification: WT\_constraint\_derived.\par

Default: ; differs from default: .\par

\textbf{Equation role.} Full conserved state and regulatory initial value at t=0 for the finite duration protocol; all amounts and both volumes evolve subsequently.\par

\textbf{Constraint.} Solved, accepted WT resting state used as the shared initial condition for production genotype trajectories; not the constructor seed and not an independent measurement.\par

\textbf{Limit of justification.} An accepted physiological point does not identify this setting or supply a numerical uncertainty interval. No knockout fit or new measurement is performed in this audit.\par

\textbf{Logical domain.} Positive for the retained law; no finite upper physiological bound established.\par

\textbf{Uncertainty status.} conditional\_root\_value\_no\_experimental\_interval\par

\textbf{Dependence of claims.} Numerical state, current and secretion claims are conditional on this active setting; exact stoichiometric identities are independent of its numerical value.\par

\textbf{Discriminating measurement.} Measure this quantity or its identifiable combination in the same preparation and stimulus protocol, with independent WT data; determine physiological ranges before making global robustness claims.\par

\textbf{Source.} \nolinkurl{src/modern_full_model/model.py:559}. No independently verified primary measurement attached to this numerical setting; model lineage is the source.\par}

\subsubsection*{Record 125: na l fmol}
{\small \nolinkurl{frozen_initial.na_l_fmol}\par

Value: \texttt{13.12776866822541}; units: fmol; active: Yes; classification: WT\_constraint\_derived.\par

Default: ; differs from default: .\par

\textbf{Equation role.} Full conserved state and regulatory initial value at t=0 for the finite duration protocol; all amounts and both volumes evolve subsequently.\par

\textbf{Constraint.} Solved, accepted WT resting state used as the shared initial condition for production genotype trajectories; not the constructor seed and not an independent measurement.\par

\textbf{Limit of justification.} An accepted physiological point does not identify this setting or supply a numerical uncertainty interval. No knockout fit or new measurement is performed in this audit.\par

\textbf{Logical domain.} Positive for the retained law; no finite upper physiological bound established.\par

\textbf{Uncertainty status.} conditional\_root\_value\_no\_experimental\_interval\par

\textbf{Dependence of claims.} Numerical state, current and secretion claims are conditional on this active setting; exact stoichiometric identities are independent of its numerical value.\par

\textbf{Discriminating measurement.} Measure this quantity or its identifiable combination in the same preparation and stimulus protocol, with independent WT data; determine physiological ranges before making global robustness claims.\par

\textbf{Source.} \nolinkurl{src/modern_full_model/model.py:559}. No independently verified primary measurement attached to this numerical setting; model lineage is the source.\par}

\subsubsection*{Record 126: k l fmol}
{\small \nolinkurl{frozen_initial.k_l_fmol}\par

Value: \texttt{2.718972608990927}; units: fmol; active: Yes; classification: WT\_constraint\_derived.\par

Default: ; differs from default: .\par

\textbf{Equation role.} Full conserved state and regulatory initial value at t=0 for the finite duration protocol; all amounts and both volumes evolve subsequently.\par

\textbf{Constraint.} Solved, accepted WT resting state used as the shared initial condition for production genotype trajectories; not the constructor seed and not an independent measurement.\par

\textbf{Limit of justification.} An accepted physiological point does not identify this setting or supply a numerical uncertainty interval. No knockout fit or new measurement is performed in this audit.\par

\textbf{Logical domain.} Positive for the retained law; no finite upper physiological bound established.\par

\textbf{Uncertainty status.} conditional\_root\_value\_no\_experimental\_interval\par

\textbf{Dependence of claims.} Numerical state, current and secretion claims are conditional on this active setting; exact stoichiometric identities are independent of its numerical value.\par

\textbf{Discriminating measurement.} Measure this quantity or its identifiable combination in the same preparation and stimulus protocol, with independent WT data; determine physiological ranges before making global robustness claims.\par

\textbf{Source.} \nolinkurl{src/modern_full_model/model.py:559}. No independently verified primary measurement attached to this numerical setting; model lineage is the source.\par}

\subsubsection*{Record 127: cl l fmol}
{\small \nolinkurl{frozen_initial.cl_l_fmol}\par

Value: \texttt{15.37868785453798}; units: fmol; active: Yes; classification: WT\_constraint\_derived.\par

Default: ; differs from default: .\par

\textbf{Equation role.} Full conserved state and regulatory initial value at t=0 for the finite duration protocol; all amounts and both volumes evolve subsequently.\par

\textbf{Constraint.} Solved, accepted WT resting state used as the shared initial condition for production genotype trajectories; not the constructor seed and not an independent measurement.\par

\textbf{Limit of justification.} An accepted physiological point does not identify this setting or supply a numerical uncertainty interval. No knockout fit or new measurement is performed in this audit.\par

\textbf{Logical domain.} Positive for the retained law; no finite upper physiological bound established.\par

\textbf{Uncertainty status.} conditional\_root\_value\_no\_experimental\_interval\par

\textbf{Dependence of claims.} Numerical state, current and secretion claims are conditional on this active setting; exact stoichiometric identities are independent of its numerical value.\par

\textbf{Discriminating measurement.} Measure this quantity or its identifiable combination in the same preparation and stimulus protocol, with independent WT data; determine physiological ranges before making global robustness claims.\par

\textbf{Source.} \nolinkurl{src/modern_full_model/model.py:559}. No independently verified primary measurement attached to this numerical setting; model lineage is the source.\par}

\subsubsection*{Record 128: tic l fmol}
{\small \nolinkurl{frozen_initial.tic_l_fmol}\par

Value: \texttt{0.5402623164056106}; units: fmol; active: Yes; classification: WT\_constraint\_derived.\par

Default: ; differs from default: .\par

\textbf{Equation role.} Full conserved state and regulatory initial value at t=0 for the finite duration protocol; all amounts and both volumes evolve subsequently.\par

\textbf{Constraint.} Solved, accepted WT resting state used as the shared initial condition for production genotype trajectories; not the constructor seed and not an independent measurement.\par

\textbf{Limit of justification.} An accepted physiological point does not identify this setting or supply a numerical uncertainty interval. No knockout fit or new measurement is performed in this audit.\par

\textbf{Logical domain.} Positive for the retained law; no finite upper physiological bound established.\par

\textbf{Uncertainty status.} conditional\_root\_value\_no\_experimental\_interval\par

\textbf{Dependence of claims.} Numerical state, current and secretion claims are conditional on this active setting; exact stoichiometric identities are independent of its numerical value.\par

\textbf{Discriminating measurement.} Measure this quantity or its identifiable combination in the same preparation and stimulus protocol, with independent WT data; determine physiological ranges before making global robustness claims.\par

\textbf{Source.} \nolinkurl{src/modern_full_model/model.py:559}. No independently verified primary measurement attached to this numerical setting; model lineage is the source.\par}

\subsubsection*{Record 129: alk l fmol}
{\small \nolinkurl{frozen_initial.alk_l_fmol}\par

Value: \texttt{0.4680534226783657}; units: fmol; active: Yes; classification: WT\_constraint\_derived.\par

Default: ; differs from default: .\par

\textbf{Equation role.} Full conserved state and regulatory initial value at t=0 for the finite duration protocol; all amounts and both volumes evolve subsequently.\par

\textbf{Constraint.} Solved, accepted WT resting state used as the shared initial condition for production genotype trajectories; not the constructor seed and not an independent measurement.\par

\textbf{Limit of justification.} An accepted physiological point does not identify this setting or supply a numerical uncertainty interval. No knockout fit or new measurement is performed in this audit.\par

\textbf{Logical domain.} Positive for the retained law; no finite upper physiological bound established.\par

\textbf{Uncertainty status.} conditional\_root\_value\_no\_experimental\_interval\par

\textbf{Dependence of claims.} Numerical state, current and secretion claims are conditional on this active setting; exact stoichiometric identities are independent of its numerical value.\par

\textbf{Discriminating measurement.} Measure this quantity or its identifiable combination in the same preparation and stimulus protocol, with independent WT data; determine physiological ranges before making global robustness claims.\par

\textbf{Source.} \nolinkurl{src/modern_full_model/model.py:559}. No independently verified primary measurement attached to this numerical setting; model lineage is the source.\par}

\subsubsection*{Record 130: volume l pL}
{\small \nolinkurl{frozen_initial.volume_l_pL}\par

Value: \texttt{0.1053884073536874}; units: pL; active: Yes; classification: WT\_constraint\_derived.\par

Default: ; differs from default: .\par

\textbf{Equation role.} Full conserved state and regulatory initial value at t=0 for the finite duration protocol; all amounts and both volumes evolve subsequently.\par

\textbf{Constraint.} Solved, accepted WT resting state used as the shared initial condition for production genotype trajectories; not the constructor seed and not an independent measurement.\par

\textbf{Limit of justification.} An accepted physiological point does not identify this setting or supply a numerical uncertainty interval. No knockout fit or new measurement is performed in this audit.\par

\textbf{Logical domain.} Positive for the retained law; no finite upper physiological bound established.\par

\textbf{Uncertainty status.} conditional\_root\_value\_no\_experimental\_interval\par

\textbf{Dependence of claims.} Numerical state, current and secretion claims are conditional on this active setting; exact stoichiometric identities are independent of its numerical value.\par

\textbf{Discriminating measurement.} Measure this quantity or its identifiable combination in the same preparation and stimulus protocol, with independent WT data; determine physiological ranges before making global robustness claims.\par

\textbf{Source.} \nolinkurl{src/modern_full_model/model.py:559}. No independently verified primary measurement attached to this numerical setting; model lineage is the source.\par}

\subsubsection*{Record 131: ae4 effective activation fraction}
{\small \nolinkurl{frozen_initial.ae4_effective_activation_fraction}\par

Value: \texttt{0}; units: dimensionless; active: Yes; classification: WT\_constraint\_derived.\par

Default: ; differs from default: .\par

\textbf{Equation role.} Full conserved state and regulatory initial value at t=0 for the finite duration protocol; all amounts and both volumes evolve subsequently.\par

\textbf{Constraint.} Solved, accepted WT resting state used as the shared initial condition for production genotype trajectories; not the constructor seed and not an independent measurement.\par

\textbf{Limit of justification.} An accepted physiological point does not identify this setting or supply a numerical uncertainty interval. No knockout fit or new measurement is performed in this audit.\par

\textbf{Logical domain.} [0,1] by fractional definition only; not a measured uncertainty interval.\par

\textbf{Uncertainty status.} conditional\_root\_value\_no\_experimental\_interval\par

\textbf{Dependence of claims.} Numerical state, current and secretion claims are conditional on this active setting; exact stoichiometric identities are independent of its numerical value.\par

\textbf{Discriminating measurement.} Measure this quantity or its identifiable combination in the same preparation and stimulus protocol, with independent WT data; determine physiological ranges before making global robustness claims.\par

\textbf{Source.} \nolinkurl{src/modern_full_model/model.py:559}. No independently verified primary measurement attached to this numerical setting; model lineage is the source.\par}

\subsubsection*{Record 132: name}
{\small \nolinkurl{genotype.name}\par

Value: WT; units: genotype label; active: No; classification: protocol\_setting.\par

Default: ; differs from default: .\par

\textbf{Equation role.} Externally passed expression multipliers scale the corresponding transport flux; zero AE4 expression is an exact deletion of its source.\par

\textbf{Constraint.} Externally passed WT expression scale, equal to one for all retained transporters; not a measured abundance. Task 40 alters only the AE4 expression scale.\par

\textbf{Limit of justification.} An accepted physiological point does not identify this setting or supply a numerical uncertainty interval. No knockout fit or new measurement is performed in this audit.\par

\textbf{Logical domain.} Selected discrete setting or stated numerical reference.\par

\textbf{Uncertainty status.} genuinely\_unknown\_for\_physiological\_transfer\par

\textbf{Dependence of claims.} No independent parent RHS sensitivity. Used only for the explicitly stated metadata, initial reference, design lineage or inverse family role.\par

\textbf{Discriminating measurement.} Measure this quantity or its identifiable combination in the same preparation and stimulus protocol, with independent WT data; determine physiological ranges before making global robustness claims.\par

\textbf{Source.} \nolinkurl{src/modern_full_model/model.py:93}. No independently verified primary measurement attached to this numerical setting; model lineage is the source.\par}

\subsubsection*{Record 133: ae4 expression}
{\small \nolinkurl{genotype.ae4_expression}\par

Value: \texttt{1}; units: dimensionless expression scale; active: Yes; classification: protocol\_setting.\par

Default: ; differs from default: .\par

\textbf{Equation role.} Externally passed expression multipliers scale the corresponding transport flux; zero AE4 expression is an exact deletion of its source.\par

\textbf{Constraint.} Externally passed WT expression scale, equal to one for all retained transporters; not a measured abundance. Task 40 alters only the AE4 expression scale.\par

\textbf{Limit of justification.} An accepted physiological point does not identify this setting or supply a numerical uncertainty interval. No knockout fit or new measurement is performed in this audit.\par

\textbf{Logical domain.} Nonnegative expression scale; values 1, 0.05 and 0 are specified interventions, not a fitted interval.\par

\textbf{Uncertainty status.} genuinely\_unknown\_for\_physiological\_transfer\par

\textbf{Dependence of claims.} Numerical state, current and secretion claims are conditional on this active setting; exact stoichiometric identities are independent of its numerical value.\par

\textbf{Discriminating measurement.} Measure this quantity or its identifiable combination in the same preparation and stimulus protocol, with independent WT data; determine physiological ranges before making global robustness claims.\par

\textbf{Source.} \nolinkurl{src/modern_full_model/model.py:93}. No independently verified primary measurement attached to this numerical setting; model lineage is the source.\par}

\subsubsection*{Record 134: ae2 expression}
{\small \nolinkurl{genotype.ae2_expression}\par

Value: \texttt{1}; units: dimensionless expression scale; active: Yes; classification: protocol\_setting.\par

Default: ; differs from default: .\par

\textbf{Equation role.} Externally passed expression multipliers scale the corresponding transport flux; zero AE4 expression is an exact deletion of its source.\par

\textbf{Constraint.} Externally passed WT expression scale, equal to one for all retained transporters; not a measured abundance. Task 40 alters only the AE4 expression scale.\par

\textbf{Limit of justification.} An accepted physiological point does not identify this setting or supply a numerical uncertainty interval. No knockout fit or new measurement is performed in this audit.\par

\textbf{Logical domain.} Nonnegative expression scale; values 1, 0.05 and 0 are specified interventions, not a fitted interval.\par

\textbf{Uncertainty status.} genuinely\_unknown\_for\_physiological\_transfer\par

\textbf{Dependence of claims.} Numerical state, current and secretion claims are conditional on this active setting; exact stoichiometric identities are independent of its numerical value.\par

\textbf{Discriminating measurement.} Measure this quantity or its identifiable combination in the same preparation and stimulus protocol, with independent WT data; determine physiological ranges before making global robustness claims.\par

\textbf{Source.} \nolinkurl{src/modern_full_model/model.py:93}. No independently verified primary measurement attached to this numerical setting; model lineage is the source.\par}

\subsubsection*{Record 135: nkcc1 expression}
{\small \nolinkurl{genotype.nkcc1_expression}\par

Value: \texttt{1}; units: dimensionless expression scale; active: Yes; classification: protocol\_setting.\par

Default: ; differs from default: .\par

\textbf{Equation role.} Externally passed expression multipliers scale the corresponding transport flux; zero AE4 expression is an exact deletion of its source.\par

\textbf{Constraint.} Externally passed WT expression scale, equal to one for all retained transporters; not a measured abundance. Task 40 alters only the AE4 expression scale.\par

\textbf{Limit of justification.} An accepted physiological point does not identify this setting or supply a numerical uncertainty interval. No knockout fit or new measurement is performed in this audit.\par

\textbf{Logical domain.} Nonnegative expression scale; values 1, 0.05 and 0 are specified interventions, not a fitted interval.\par

\textbf{Uncertainty status.} genuinely\_unknown\_for\_physiological\_transfer\par

\textbf{Dependence of claims.} Numerical state, current and secretion claims are conditional on this active setting; exact stoichiometric identities are independent of its numerical value.\par

\textbf{Discriminating measurement.} Measure this quantity or its identifiable combination in the same preparation and stimulus protocol, with independent WT data; determine physiological ranges before making global robustness claims.\par

\textbf{Source.} \nolinkurl{src/modern_full_model/model.py:93}. No independently verified primary measurement attached to this numerical setting; model lineage is the source.\par}

\subsubsection*{Record 136: nhe1 expression}
{\small \nolinkurl{genotype.nhe1_expression}\par

Value: \texttt{1}; units: dimensionless expression scale; active: Yes; classification: protocol\_setting.\par

Default: ; differs from default: .\par

\textbf{Equation role.} Externally passed expression multipliers scale the corresponding transport flux; zero AE4 expression is an exact deletion of its source.\par

\textbf{Constraint.} Externally passed WT expression scale, equal to one for all retained transporters; not a measured abundance. Task 40 alters only the AE4 expression scale.\par

\textbf{Limit of justification.} An accepted physiological point does not identify this setting or supply a numerical uncertainty interval. No knockout fit or new measurement is performed in this audit.\par

\textbf{Logical domain.} Nonnegative expression scale; values 1, 0.05 and 0 are specified interventions, not a fitted interval.\par

\textbf{Uncertainty status.} genuinely\_unknown\_for\_physiological\_transfer\par

\textbf{Dependence of claims.} Numerical state, current and secretion claims are conditional on this active setting; exact stoichiometric identities are independent of its numerical value.\par

\textbf{Discriminating measurement.} Measure this quantity or its identifiable combination in the same preparation and stimulus protocol, with independent WT data; determine physiological ranges before making global robustness claims.\par

\textbf{Source.} \nolinkurl{src/modern_full_model/model.py:93}. No independently verified primary measurement attached to this numerical setting; model lineage is the source.\par}

\subsubsection*{Record 137: ae4 expression cases}
{\small \nolinkurl{protocol.ae4_expression_cases}\par

Value: [1.0, 0.05, 0.0]; units: dimensionless expression scales; active: Yes; classification: protocol\_setting.\par

Default: ; differs from default: .\par

\textbf{Equation role.} Task40 calls the same model with e4=1,0.05,0; eNKCC=eNHE=eAE2=1 and all other parameters fixed.\par

\textbf{Constraint.} Predeclared WT, 5 percent and null interventions; only the AE4 expression field differs. These are protocol conditions, not estimated abundance uncertainty.\par

\textbf{Limit of justification.} An accepted physiological point does not identify this setting or supply a numerical uncertainty interval. No knockout fit or new measurement is performed in this audit.\par

\textbf{Logical domain.} Selected discrete setting or stated numerical reference.\par

\textbf{Uncertainty status.} genuinely\_unknown\_for\_physiological\_transfer\par

\textbf{Dependence of claims.} Numerical state, current and secretion claims are conditional on this active setting; exact stoichiometric identities are independent of its numerical value.\par

\textbf{Discriminating measurement.} Measure this quantity or its identifiable combination in the same preparation and stimulus protocol, with independent WT data; determine physiological ranges before making global robustness claims.\par

\textbf{Source.} \nolinkurl{analysis/40_ae4_equal_cation_routing/run_case.py:78}. No independently verified primary measurement attached to this numerical setting; model lineage is the source.\par}
